%% ============================================================
%%  Série Avaliação na Prática — FGV CLEAR
%%  Guia de Coleta de Dados Primários
%% ============================================================

% !TEX program = pdflatex
% !TEX encoding = UTF-8 Unicode

\documentclass[12pt, a4paper, openany]{book}

%% ── Codificação e Idioma ────────────────────────────────────
\usepackage[utf8]{inputenc}
\usepackage[T1]{fontenc}
\usepackage[brazil]{babel}

%% ── Fonte ───────────────────────────────────────────────────
\usepackage{tgpagella}

%% ── Matemática ──────────────────────────────────────────────
\usepackage{amsmath, amssymb, amsfonts, dsfont, amsthm}

%% ── Figuras e Tabelas ───────────────────────────────────────
\usepackage{graphicx, float}
\usepackage[labelfont=bf, font=small, justification=justified]{caption}
\usepackage{subcaption}
\usepackage{booktabs, multirow, multicol, array}
\usepackage{longtable}
\usepackage{tabularx}

%% ── Layout ──────────────────────────────────────────────────
\usepackage{geometry}
\geometry{a4paper, left=2.5cm, right=2.5cm, top=2.5cm, bottom=2.5cm}
\usepackage{setspace}
\onehalfspacing
\usepackage{fancyhdr}
\usepackage{emptypage}
\usepackage{enumitem}

%% ── Cores (paleta CLEAR) ────────────────────────────────────
\usepackage{xcolor}
\definecolor{clearblue}{HTML}{1B3A6E}
\definecolor{clearlightblue}{HTML}{5CA4A9}
\definecolor{cleardarkblue}{HTML}{003F72}
\definecolor{cleargray}{HTML}{6E6E6E}
\definecolor{salmonbg}{RGB}{252,235,230}
\definecolor{salmonframe}{RGB}{230,160,140}
\definecolor{bluebg}{RGB}{235,245,251}
\definecolor{blueframe}{RGB}{58,125,181}
\definecolor{orangebg}{RGB}{254,245,231}
\definecolor{orangeframe}{RGB}{230,126,34}
\definecolor{codebg}{RGB}{248,248,248}
\definecolor{codeframe}{RGB}{200,200,200}
\definecolor{redbg}{RGB}{250,230,225}
\definecolor{redframe}{RGB}{180,40,40}
\definecolor{greenbg}{RGB}{225,242,225}
\definecolor{greenframe}{RGB}{60,140,80}

%% ── Hyperlinks e Bibliografia ───────────────────────────────
\usepackage{hyperref}
\hypersetup{
  colorlinks=true,
  linkcolor=clearlightblue,
  citecolor=clearlightblue,
  urlcolor=clearlightblue,
  pdfauthor={FGV CLEAR},
  pdftitle={Guia de Coleta de Dados Primários}
}
\usepackage{natbib}

%% ── TColorbox ───────────────────────────────────────────────
\usepackage[many]{tcolorbox}
\tcbuselibrary{breakable, skins}

%% ── Código ──────────────────────────────────────────────────
\usepackage{listings}

%% ── Cabeçalhos e Rodapés ────────────────────────────────────
\pagestyle{fancy}
\fancyhf{}
\fancyhead[R]{\small\color{clearblue}\textit{Guia CLEAR $|$ Coleta de Dados Primários}}
\fancyfoot[C]{\thepage}
\renewcommand{\headrulewidth}{0.4pt}
\renewcommand{\headrule}{\hbox to\headwidth{\color{clearblue}\leaders\hrule height \headrulewidth\hfill}}

%% ── Títulos de Seção ────────────────────────────────────────
\usepackage{titlesec}
\titleformat{\chapter}[display]
  {\normalfont\huge\bfseries\color{clearblue}}
  {\chaptertitlename\ \thechapter}{20pt}{\Huge}
\titleformat{\section}
  {\normalfont\Large\bfseries\color{clearblue}}{\thesection}{1em}{}
\titleformat{\subsection}
  {\normalfont\large\bfseries\color{cleardarkblue}}{\thesubsection}{1em}{}
\titleformat{\subsubsection}
  {\normalfont\normalsize\bfseries\color{cleardarkblue}}{\thesubsubsection}{1em}{}

%% ── Caixas Temáticas ────────────────────────────────────────
\newtcolorbox{destaquebox}[1][Destaque]{
  enhanced, breakable,
  colback=salmonbg, colframe=salmonframe, coltitle=black, fonttitle=\bfseries,
  title=#1, boxrule=0.5pt, arc=2pt,
  left=6pt, right=6pt, top=4pt, bottom=4pt
}
\newtcolorbox{notabox}[1][Nota]{
  enhanced, breakable,
  colback=bluebg, colframe=blueframe, coltitle=white, fonttitle=\bfseries,
  title=#1, boxrule=0.5pt, arc=2pt,
  left=6pt, right=6pt, top=4pt, bottom=4pt
}
\newtcolorbox{conexaobox}[1][Conex\~ao com a s\'erie]{
  enhanced, breakable,
  colback=orangebg, colframe=orangeframe, coltitle=black, fonttitle=\bfseries,
  title=#1, boxrule=0.5pt, arc=2pt,
  left=6pt, right=6pt, top=4pt, bottom=4pt
}
\newtcolorbox{exemplobox}[1][Exemplo]{
  enhanced, breakable,
  colback=greenbg, colframe=greenframe, coltitle=white, fonttitle=\bfseries,
  title=#1, boxrule=0.5pt, arc=2pt,
  left=6pt, right=6pt, top=4pt, bottom=4pt
}
\newtcolorbox{atencaobox}[1][Aten\c{c}\~ao]{
  enhanced, breakable,
  colback=redbg, colframe=redframe, coltitle=white, fonttitle=\bfseries,
  title=#1, boxrule=0.5pt, arc=2pt,
  left=6pt, right=6pt, top=4pt, bottom=4pt
}


%% ── Estilo de Listings (caso necessário) ────────────────────
\lstdefinestyle{Rstyle}{
  language=R,
  backgroundcolor=\color{codebg},
  frame=single, rulecolor=\color{codeframe},
  basicstyle=\small\ttfamily,
  numbers=none, breaklines=true,
  showstringspaces=false, tabsize=2,
  xleftmargin=6pt, xrightmargin=6pt,
  aboveskip=8pt, belowskip=8pt
}

%% ── Metadados de Figuras/Quadros ────────────────────────────
\newcommand{\figmeta}[3]{%
  \par\vspace{4pt}%
  \begin{minipage}{\linewidth}%
  \footnotesize%
  \textbf{Fonte:} #1%
  \ifx&#2&\else\\\textbf{Elabora\c{c}\~ao:} #2\fi%
  \ifx&#3&\else\\\textbf{Nota:} #3\fi%
  \end{minipage}%
  \vspace{6pt}%
}

%% ── Profundidade de numera\c{c}\~ao ─────────────────────────
\setcounter{secnumdepth}{2}
\setcounter{tocdepth}{2}

% ============================================================
\begin{document}

% ============================================================
% CAPA
% ============================================================
\thispagestyle{empty}
\begin{center}
\vspace*{3cm}

{\Large\bfseries\color{clearlightblue} S\'ERIE: AVALIA\c{C}\~AO NA PR\'ATICA}

\vspace{1.8cm}

{\Huge\bfseries\color{clearblue} Guia de Coleta de\\[0.3cm]
Dados Prim\'arios}

\vspace{2.2cm}

{\large FGV CLEAR $|$ FGV EESP}

\vspace{0.5cm}

{\large Maio de 2026}

\vfill
\end{center}
\newpage

% ============================================================
% FICHA TÉCNICA
% ============================================================
\thispagestyle{empty}
\vspace*{1.5cm}

{\Large\bfseries\color{clearblue} Ficha T\'ecnica}

\bigskip

\noindent\textbf{Equipe t\'ecnica:} Beatriz Burattini, Caio de Souza Castro, Pedro Molina Ogeda, Vitor Menezes.

\bigskip

\noindent\textbf{Apresenta\c{c}\~ao.} Este guia integra a s\'erie de publica\c{c}\~oes \textit{Avalia\c{c}\~ao na Pr\'atica}, desenvolvida pelo FGV CLEAR com o objetivo de ampliar o acesso a conhecimentos sobre monitoramento e avalia\c{c}\~ao com foco em pol\'iticas p\'ublicas.

\medskip

\noindent Fundado em 2015, o FGV CLEAR tem se dedicado ao fortalecimento da cultura de gest\~ao orientada por evid\^encias no Brasil e em pa\'ises lus\'ofonos. Com sede na Escola de Economia de S\~ao Paulo da Funda\c{c}\~ao Getulio Vargas (FGV EESP), o FGV CLEAR atua como centro regional da Iniciativa CLEAR (\textit{Centers for Learning on Evaluation and Results}).

\medskip

\noindent Saiba mais sobre o FGV CLEAR e acesse outras publica\c{c}\~oes em: \textbf{www.fgvclear.org}

\vfill

\noindent{\small\color{cleargray}%
Os textos desta publica\c{c}\~ao s\~ao de responsabilidade exclusiva dos autores. As opini\~oes expressas n\~ao refletem necessariamente a posi\c{c}\~ao das institui\c{c}\~oes parceiras.}

\newpage

% ============================================================
% SUMÁRIO
% ============================================================
% ─── SOBRE ESTE GUIA ─────────────────────────────────────────────────────────
\thispagestyle{empty}
\section*{Sobre este guia}
Nem toda pergunta de avaliação tem resposta numa base existente. Saber
se um curso de qualificação melhorou a autopercepção dos participantes
sobre suas chances de conseguir emprego; se um programa de mediação
escolar mudou a forma como alunos lidam com conflitos; se um treinamento
de gestores afetou a forma como conduzem reuniões; se moradores de uma
região se sentem mais seguros depois da reforma da praça do bairro: nada
disso está no CAGED, na PNAD, no Datasus, nem em qualquer registro
administrativo. Para perguntas sobre percepção, comportamento, prática,
opinião ou processo, os dados existentes raramente bastam. O avaliador
precisa gerá-los.

A coleta primária é o processo de gerar esses dados em campo. Cada
decisão envolvida afeta o que sai do outro lado, e a sequência de
decisões importa tanto quanto a sofisticação da análise que vem depois.
Uma análise impecável de dados ruins ainda é uma análise de dados ruins;
nenhum método consegue extrair de uma amostra o que ela não tem.

A primeira decisão é a metodologia. Survey presencial, em domicílio ou em
local específico, é caro mas costuma ter taxa de resposta alta e permite
controle do entrevistador sobre como cada pergunta é interpretada. Survey
por telefone é mais barato e mais rápido, mas filtra populações sem linha
estável e produz amostras enviesadas para certos perfis. Survey online
(cada vez mais usado) tem custo marginal próximo de zero e permite coleta
em larga escala, mas amostra de fato quem tem acesso à internet e
disposição para responder, o que pode excluir partes importantes da
população de interesse. Cada modo de coleta vem com seu viés
característico, e a escolha está longe de ser independente da pergunta
sendo feita.

A segunda é o plano amostral. Pode-se sortear participantes da população
inteira com probabilidade uniforme (amostra aleatória simples), separar
a população em estratos relevantes e sortear dentro de cada um para
garantir representação (estratificada), ou amostrar conglomerados
inteiros para reduzir custo logístico (cluster). A escolha entre essas
estratégias afeta tanto o custo do trabalho quanto a precisão da
estimativa final, e o cálculo do tamanho de amostra necessário depende
dela. Amostrar mal é a forma mais comum de produzir dados que parecem
informativos sobre uma população mas só descrevem o subgrupo que
respondeu, que costuma ser o mais fácil de encontrar e também o mais
diferente da população que se queria estudar.

A terceira é a redação dos instrumentos. Como cada pergunta é formulada
determina o que o respondente entende dela, o que ele decide responder, o
que ele decide omitir, e em que escala ele organiza sua resposta. Uma
pergunta ambígua produz respostas que não são comparáveis entre
indivíduos. Uma pergunta enviesada (que sugere a resposta esperada)
produz a resposta esperada. Uma escala mal calibrada (cinco pontos onde
os extremos quase nunca são usados) descarta variação útil. Uma pergunta
sensível mal posicionada no questionário (perguntada cedo demais, ou logo
depois de uma pergunta que predispõe a determinada resposta) contamina o
que se mediria. Esses problemas são difíceis de detectar depois de
coletados os dados e impossíveis de consertar com mais análise. A única
defesa é desenhar bem antes, com piloto, revisão por pessoas com perfil
próximo ao da amostra, e cuidado na ordem e no enquadramento de cada
item.

A quarta é a aplicação em campo. Treinamento dos entrevistadores,
supervisão durante a coleta, controle de qualidade conforme os dados
entram, gestão de não-resposta. A qualidade do dado final depende muito
mais desses processos, frequentemente esquecidos na fase de planejamento,
do que da metodologia escolhida ou do instrumento desenhado.

Por trás dessas quatro decisões há uma quinta que cobre todas: ética em
pesquisa com seres humanos. Consentimento informado, confidencialidade,
justificativa do incômodo causado ao respondente, devolutiva apropriada
quando relevante. No Brasil, projetos de pesquisa com pessoas precisam
passar por um Comitê de Ética em Pesquisa (CEP), processo que tem suas
peculiaridades e que o avaliador precisa conhecer antes de chegar ao
momento de aplicar o questionário.

Este guia trata dessas decisões na sequência em que costumam aparecer num
projeto real: metodologia, amostra, ética, elaboração de instrumentos, e
aplicação. É indicado para quem precisa montar uma pesquisa do zero, para
diagnóstico, monitoramento ou avaliação.
\newpage

% ─── COMO LER ESTE GUIA ──────────────────────────────────────────────────────
\thispagestyle{empty}
\section*{Como ler este guia}
Três tipos de caixa estruturam a leitura, identificados por cor:
\vspace{0.3cm}
\begin{notabox}[Caixa azul --- pontos-chave]
Aparece ao fim de cada capítulo, com sínteses e recomendações para quem
precisa decidir.
\end{notabox}
\begin{exemplobox}[Caixa verde --- exemplo aplicado]
Mostra como as ideias do capítulo se traduzem em situações de política
pública. Os casos são fictícios, mas calibrados por padrões recorrentes
na literatura e na prática de gestão.
\end{exemplobox}
\begin{atencaobox}[Caixa vermelha --- cuidados essenciais]
Sinaliza armadilhas comuns --- pontos em que uma escolha mal fundamentada
compromete toda a análise.
\end{atencaobox}
\vspace{0.5cm}
\noindent O guia faz referência recorrente ao \textit{Guia CLEAR de
Monitoramento e Avaliação de Políticas Públicas: do diagnóstico à decisão}
\citep{Lima2025} --- doravante \textbf{Guia CLEAR de M\&A} ---, que trata
em profundidade as etapas do ciclo da política pública adotado na série.
Quando o texto menciona Teoria da Mudança, plano de monitoramento ou
tipologia de avaliações, é nesse guia que o leitor encontra o
desenvolvimento completo.
\newpage

\tableofcontents
\newpage

% ============================================================
% 1. INTRODUÇÃO
% ============================================================

\chapter{Introdu\c{c}\~ao}

Decis\~oes de pol\'itica p\'ublica que se pretendem informadas dependem de algo banal e dif\'icil: informa\c{c}\~ao confi\'avel sobre o que est\'a acontecendo no campo. Registros administrativos e pesquisas j\'a existentes resolvem parte do problema --- s\~ao baratos, est\~ao prontos e cobrem dimens\~oes amplas da realidade. Mas raramente cobrem tudo. Quando a pergunta avaliativa exige informa\c{c}\~ao que essas bases n\~ao t\^em (ou n\~ao t\^em na granularidade necess\'aria, ou n\~ao t\^em atualizada), n\~ao h\'a alternativa: \'e preciso coletar.

A coleta prim\'aria \'e a produ\c{c}\~ao de informa\c{c}\~oes originais junto a quem participa, se beneficia ou implementa uma pol\'itica. \'E o canal pelo qual a avalia\c{c}\~ao acessa n\~ao apenas resultados, mas processos, percep\c{c}\~oes, barreiras e contextos --- as dimens\~oes que registros administrativos t\^ipicamente n\~ao registram. \'E tamb\'em por onde se identificam efeitos n\~ao previstos, se distinguem falhas de teoria de falhas de implementa\c{c}\~ao e se mapeiam mudan\c{c}as estruturais que indicadores convencionais n\~ao captam.

Este Guia foi escrito para quem desenvolve, executa ou supervisiona pesquisas e avalia\c{c}\~oes. O texto trata das decis\~oes metodol\'ogicas centrais, dos cuidados \'eticos e operacionais e das boas pr\'aticas que sustentam a validade e a confiabilidade das informa\c{c}\~oes produzidas em campo.

A organiza\c{c}\~ao do guia segue o fluxo de uma pesquisa aplicada. O Cap\'itulo~2 trata das metodologias de coleta, contrastando abordagens quantitativas e qualitativas. O Cap\'itulo~3 aborda a defini\c{c}\~ao da amostra e os crit\'erios de sele\c{c}\~ao de participantes. O Cap\'itulo~4 discute as considera\c{c}\~oes \'eticas --- princ\'ipios gerais, marco normativo, sele\c{c}\~ao de participantes e tratamento dos dados ---, que atravessam todo o processo. O Cap\'itulo~5 entra na elabora\c{c}\~ao de question\'arios e roteiros, com \^enfase em clareza, precis\~ao e adequa\c{c}\~ao ao objetivo. O Cap\'itulo~6 fecha com os procedimentos de campo: treinamento, acompanhamento, prote\c{c}\~ao dos dados e controle de qualidade. As Considera\c{c}\~oes finais retomam limites de generaliza\c{c}\~ao e quest\~oes operacionais que perpassam tudo o que veio antes.

% ============================================================
% 2. METODOLOGIAS DE COLETA
% ============================================================
\chapter{Metodologias de coleta}

A escolha da metodologia de coleta depende dos objetivos da pesquisa, do perfil do p\'ublico-alvo, dos recursos dispon\'iveis e do tipo de informa\c{c}\~ao buscada. \textit{Surveys} (ou question\'arios) entregam dados quantitativos compar\'aveis em amostras grandes e permitem an\'alises generaliz\'aveis. Entrevistas trazem profundidade --- contexto, l\'ogicas sociais, motivos por tr\'as do comportamento --- e por isso s\~ao a via natural quando o que se quer entender s\~ao percep\c{c}\~oes e experi\^encias.

A Tabela~\ref{tab:quando_quanti_quali} sintetiza a refer\^encia elaborada por \citet{morra2009road} para a escolha do instrumento adequado a cada tipo de pesquisa avaliativa. T\'ecnicas quantitativas e qualitativas representam paradigmas distintos de investiga\c{c}\~ao, e a melhor forma de ler suas diferen\c{c}as \'e pelo que cada uma entrega:

As \textbf{abordagens quantitativas} produzem grande volume de informa\c{c}\~oes padronizadas, viabilizando mensura\c{c}\~ao precisa e an\'alise estat\'istica de fen\^omenos observ\'aveis em popula\c{c}\~oes amplas. O foco \'e quantifica\c{c}\~ao, generaliza\c{c}\~ao e identifica\c{c}\~ao de padr\~oes e rela\c{c}\~oes de causalidade.

As \textbf{abordagens qualitativas} buscam significados, percep\c{c}\~oes, motiva\c{c}\~oes e din\^amicas sociais que est\~ao por tr\'as dos fen\^omenos, sem recorrer necessariamente \`a quantifica\c{c}\~ao. Seu valor est\'a em explorar a complexidade de processos e contextos espec\'ificos.

Em muitas pesquisas avaliativas, nenhum dos dois isoladamente d\'a conta. A \textbf{metodologia mista} \citep{creswell2014research,bryman2012social} combina os dois e captura tanto a abrang\^encia dos \textit{surveys} quanto a profundidade das entrevistas e grupos focais --- arranjo particularmente \'util quando o objetivo \'e n\~ao s\'o medir o impacto de uma interven\c{c}\~ao, mas tamb\'em entender por quais mecanismos esse impacto ocorre \citep{morra2009road}.

\begin{table}[H]
\centering
\caption{Quando utilizar abordagens quantitativas e qualitativas?}
\label{tab:quando_quanti_quali}
\begin{tabular}{p{0.55\linewidth} p{0.35\linewidth}}
\toprule
\textbf{Se voc\^e} & \textbf{Ent\~ao, use esta abordagem} \\
\midrule
Deseja conduzir an\'alises estat\'isticas & Quantitativa \\
Deseja precis\~ao na quantifica\c{c}\~ao de fen\^omenos & \\
Sabe exatamente o que precisa medir & \\
Possui interesse em cobrir informa\c{c}\~oes sobre um grupo extenso de indiv\'iduos & \\
\midrule
Deseja obter informa\c{c}\~oes com mais profundidade & Qualitativa \\
Possui interesse em compreender as narrativas sobre um fen\^omeno & \\
N\~ao necessita quantificar os resultados & \\
Ainda n\~ao possui certeza sobre o que \'e capaz de medir & \\
\bottomrule
\end{tabular}
\figmeta{Adaptado de \citet{morra2009road}.}{}{}
\end{table}

Definida a abordagem, o passo seguinte \'e o instrumento. Esta se\c{c}\~ao trata da elabora\c{c}\~ao das ferramentas de coleta, presenciais e remotas, com aten\c{c}\~ao aos pontos cr\'iticos de validade, confiabilidade e adequa\c{c}\~ao ao contexto de aplica\c{c}\~ao.

\section{Metodologias de coleta quantitativa}

Dados quantitativos s\~ao informa\c{c}\~oes num\'ericas que permitem mensurar fen\^omenos objetivamente e aplicar t\'ecnicas estat\'isticas de an\'alise, compara\c{c}\~ao e generaliza\c{c}\~ao \citep{kabir2016methods,taherdoost2021data}. Sua coleta exige estrutura\c{c}\~ao pr\'evia, normalmente a partir de amostras aleat\'orias.

Comparados a m\'etodos qualitativos, m\'etodos quantitativos costumam ser mais vi\'aveis em custo e tempo, e os resultados se prestam a padroniza\c{c}\~ao, compara\c{c}\~ao e sumariza\c{c}\~ao. Em troca, perdem em profundidade contextual e podem encontrar dificuldades operacionais durante a implementa\c{c}\~ao \citep{taherdoost2021data}.

A natureza dos dados quantitativos depende das \textbf{escalas de mensura\c{c}\~ao}, que definem como cada vari\'avel pode ser interpretada e tratada estatisticamente:

\begin{itemize}
\item \textbf{Nominal:} Categoriza dados sem uma ordem hier\'arquica (ex.: g\^enero, pa\'is de origem).
\item \textbf{Ordinal:} Categoriza dados com uma ordem ou classifica\c{c}\~ao, mas sem dist\^ancia definida entre os pontos (ex.: escala de satisfa\c{c}\~ao: ``insatisfeito'', ``neutro'', ``satisfeito'').
\item \textbf{Intervalar:} Estabelece uma ordem com dist\^ancias iguais entre os pontos, mas sem um zero absoluto (ex.: temperatura em Celsius).
\item \textbf{Raz\~ao:} Possui todas as caracter\'isticas da escala intervalar, incluindo um zero absoluto que indica aus\^encia do fen\^omeno (ex.: altura, peso, renda).
\end{itemize}

Al\'em disso, os instrumentos de coleta frequentemente utilizam escalas de mensura\c{c}\~ao, como as \textbf{escalas de avalia\c{c}\~ao} (\textit{Rating Scales}), que atribuem um valor num\'erico a categorias (ex.: ``de 1 a 5, onde 1 \'e `discordo totalmente' ''), e as \textbf{escalas de atitude}, mais complexas, que avaliam a predisposi\c{c}\~ao de indiv\'iduos perante um objeto ou ideia (ex.: escala Likert).

A principal ferramenta para coleta de dados quantitativos prim\'arios \'e o \textbf{\textit{survey} estruturado} --- um question\'ario padronizado com perguntas e ordem fixas, aplicado a uma amostra de pessoas respondentes. A padroniza\c{c}\~ao garante que todas as pessoas participantes recebam o mesmo est\'imulo, o que sustenta a confiabilidade e a comparabilidade das an\'alises estat\'isticas \citep{bryman2012social}.

A l\'ogica b\'asica do \textit{survey} \'e a mesma em todas as suas vers\~oes, mas o canal de aplica\c{c}\~ao --- presencial, telef\^onico ou online --- altera o que se ganha e o que se perde em qualidade, custo e cobertura. As tr\^es modalidades coexistem na pr\'atica e atendem a perfis diferentes de pesquisa.

\subsubsection*{1) \textit{Surveys} presenciais (face a face)}

Aplicado por entrevistadora ou entrevistador treinado, em intera\c{c}\~ao direta com a pessoa respondente, o \textit{survey} presencial \'e o formato mais tradicional da coleta em pesquisa social. Hoje o registro \'e feito majoritariamente em \textit{tablets} ou \textit{smartphones}, em \textit{softwares} de coleta assistida por computador (CAPI --- \textit{Computer-Assisted Personal Interviewing}), que reduzem erros de digita\c{c}\~ao e permitem checagens autom\'aticas de consist\^encia.

\'E o m\'etodo dominante nos levantamentos de grande escala. No Brasil, o IBGE usa entrevistas presenciais na PNAD Cont\'inua e no Censo Demogr\'afico. No plano internacional, o \textit{Living Standards Measurement Study} (LSMS) do Banco Mundial \'e o caso paradigm\'atico em pa\'ises de baixa e m\'edia renda. Pesquisas eleitorais conduzidas por institutos privados tamb\'em recorrem a entrevistas presenciais em domic\'ilio ou em pontos de fluxo, em busca de representatividade do eleitorado.

A vantagem central est\'a nas altas taxas de resposta: a intera\c{c}\~ao pessoal reduz a probabilidade de recusa. O formato tamb\'em comporta question\'arios longos e complexos, com uso de cart\~oes de apoio, escalas visuais e roteiros detalhados, e permite \`a entrevistadora esclarecer d\'uvidas em tempo real --- o que aumenta a qualidade dos dados e diminui erros de compreens\~ao.

O custo, em troca, \'e alto. A coleta exige log\'istica robusta (transporte, deslocamentos, remunera\c{c}\~ao por tempo em campo) e treinamento intensivo das equipes em t\'ecnicas de abordagem e protocolos \'eticos. H\'a ainda o risco de vi\'es da pessoa entrevistadora: presen\c{c}a f\'isica, tom de voz e julgamentos impl\'icitos podem influenciar respostas de forma sistem\'atica \citep{neuman2012designing,bryman2012social}.

A implementa\c{c}\~ao envolve etapas encadeadas. O question\'ario precisa ser programado em plataforma compat\'ivel com dispositivos m\'oveis, com l\'ogica de saltos (perguntas condicionadas a respostas pr\'evias), filtros e checagens. A equipe deve ser treinada nos conte\'udos e em t\'ecnicas de entrevista. O planejamento log\'istico inclui mapeamento das \'areas de coleta, rotas de visita e listas de domic\'ilios. Durante o campo, recomenda-se supervis\~ao constante --- presencial ou via monitoramento remoto dos dispositivos. Por fim, os dados precisam ser sincronizados diariamente com servidor central, onde passam por verifica\c{c}\~oes autom\'aticas que permitem identificar e corrigir problemas ainda durante a coleta.

\subsubsection*{2) \textit{Surveys} por telefone (CATI)}

Os \textit{surveys} telef\^onicos ocupam um meio-termo de custo entre o presencial e o online. Ganharam tra\c{c}\~ao nos anos 1980 e 1990, com a difus\~ao do telefone fixo, e hoje s\~ao operados majoritariamente por sistemas CATI (\textit{Computer-Assisted Telephone Interviewing}), que automatizam a aplica\c{c}\~ao, reduzem erros de registro e facilitam supervis\~ao e armazenagem.

S\~ao r\'apidos e cobrem grandes \'areas geogr\'aficas sem deslocamento f\'isico. O anonimato relativo da voz tende a estimular respostas mais sinceras em quest\~oes delicadas (renda, comportamento sexual, opini\~ao pol\'itica). E a supervis\~ao \'e operacionalmente simples: como as entrevistadoras trabalham em centrais de atendimento, supervisoras podem acompanhar entrevistas em tempo real e intervir quando preciso. \'E o m\'etodo dominante em pesquisas eleitorais, em inqu\'eritos de sa\'ude conduzidos por minist\'erios e secretarias (monitoramento de sintomas em epidemias, por exemplo) e em pesquisas de satisfa\c{c}\~ao de bancos, telefonia e empresas de servi\c{c}os.

As limita\c{c}\~oes s\~ao tr\^es. \textit{Cobertura}: o m\'etodo exclui quem n\~ao tem telefone ou cujo n\'umero n\~ao consta dos cadastros usados na amostragem --- no Brasil, o uso de telefones fixos caiu fortemente nas duas \'ultimas d\'ecadas, agravando o problema. \textit{Dura\c{c}\~ao}: sem apoio visual, o question\'ario precisa ser curto e direto, o que limita a profundidade. \textit{Recusa}: desligar o telefone \'e muito mais f\'acil do que recusar algu\'em frente a frente, e as taxas refletem isso. H\'a ainda um vi\'es ligado ao hor\'ario da liga\c{c}\~ao: chamadas diurnas alcan\c{c}am desproporcionalmente aposentados, desempregados e quem trabalha de casa; chamadas noturnas, ao contr\'ario, encontram mais trabalhadores ativos, mas tendem a ser percebidas como intrusivas.

A implementa\c{c}\~ao come\c{c}a pela base de n\'umeros: \'e preciso constru\'i-la ou adquiri-la com representatividade da popula\c{c}\~ao-alvo, muitas vezes em parceria com operadoras ou via discagem aleat\'oria (\textit{Random Digit Dialing} --- RDD). O question\'ario precisa ser adaptado ao formato telef\^onico --- perguntas curtas, claras, de f\'acil compreens\~ao oral. O treinamento da equipe foca em comunica\c{c}\~ao por voz, empatia e manejo de obje\c{c}\~oes, j\'a que n\~ao h\'a recursos visuais. Protocolos de liga\c{c}\~ao definem hor\'arios de contato, n\'umero m\'aximo de tentativas e intervalos entre chamadas. Durante a coleta, supervisoras acompanham entrevistas em tempo real e monitoram indicadores de produtividade (tempo m\'edio de entrevista, taxa de sucesso por liga\c{c}\~ao).

\subsubsection*{3) \textit{Surveys} autoadministrados online (web/e-mail)}

No formato autoadministrado online, a pessoa respondente acessa o question\'ario por conta pr\'opria e registra suas respostas --- normalmente via \textit{link} por e-mail, redes sociais, \textit{website} ou aplicativo m\'ovel. \'E o canal de mais r\'apido crescimento, impulsionado pela difus\~ao da internet e dos \textit{smartphones}, que tornam vi\'avel alcan\c{c}ar milhares de pessoas em poucos dias e com custo marginal pr\'oximo de zero.

O custo \'e o atrativo mais imediato, junto com a rapidez. A pessoa respondente tamb\'em ganha em conveni\^encia: completa o question\'ario no seu pr\'oprio ritmo, em hor\'arios que escolhe. As plataformas suportam l\'ogica automatizada --- randomiza\c{c}\~ao de perguntas, saltos condicionais, embaralhamento de alternativas --- que reduzem vieses de ordem e enriquecem a experi\^encia. \'E o formato dominante em pesquisas de satisfa\c{c}\~ao p\'os-compra, avalia\c{c}\~oes de cursos e docentes em universidades, e foi central em situa\c{c}\~oes emergenciais como a pandemia de Covid-19, em que inqu\'eritos r\'apidos sobre sintomas, ades\~ao a medidas de preven\c{c}\~ao e vacina\c{c}\~ao foram operados em larga escala por meios eletr\^onicos.

Em contrapartida, o m\'etodo carrega vi\'es de cobertura severo em pa\'ises marcados por desigualdade digital: popula\c{c}\~oes com acesso limitado \`a internet ou baixa familiaridade digital tendem a ficar de fora, comprometendo a representatividade. A taxa de resposta tamb\'em \'e baixa, em parte pelo excesso de convites que circulam online e pela tend\^encia das pessoas usu\'arias a ignorar e-mails de pesquisa. H\'a ainda pouco controle sobre o ambiente em que a resposta \'e dada --- distra\c{c}\~ao, pressa, preenchimento aleat\'orio --- e, em alguns casos, risco de respostas m\'ultiplas pela mesma pessoa, o que exige autentica\c{c}\~ao ou checagens de consist\^encia.

A implementa\c{c}\~ao come\c{c}a pela plataforma. A escolha precisa contemplar personaliza\c{c}\~ao do question\'ario, acompanhamento das respostas em tempo real e exporta\c{c}\~ao eficiente dos dados. O \textit{design} deve ser otimizado para celular, j\'a que esse \'e o dispositivo dominante de acesso. Os pr\'e-testes precisam simular percursos l\'ogicos, verificar randomiza\c{c}\~oes e validar todos os formatos de resposta antes do lan\c{c}amento. O recrutamento combina convites por e-mail, campanhas em redes sociais, divulga\c{c}\~ao em \textit{websites} e parcerias institucionais. Lembretes autom\'aticos e monitoramento cont\'inuo do engajamento ajudam a reduzir a n\~ao resposta e a ajustar a estrat\'egia de divulga\c{c}\~ao durante o campo.

A Tabela~\ref{tab:comparacao_surveys} resume vantagens, desvantagens e passos essenciais de cada modalidade.

\begin{table}[H]
\centering
\caption{Compara\c{c}\~ao entre os m\'etodos de \textit{survey}}
\label{tab:comparacao_surveys}
\small
\begin{tabular}{p{0.13\linewidth} p{0.27\linewidth} p{0.25\linewidth} p{0.27\linewidth}}
\toprule
\textbf{M\'etodo} & \textbf{Vantagens} & \textbf{Desvantagens} & \textbf{Passos essenciais} \\
\midrule
\textbf{Presencial} &
Altas taxas de resposta; permite question\'arios longos e complexos; possibilidade de esclarecer d\'uvidas; maior qualidade e riqueza dos dados. &
Alto custo log\'istico e de pessoal; processo mais demorado; risco de vi\'es do entrevistador. &
Programar question\'ario em dispositivos m\'oveis; treinar intensivamente entrevistadores; planejar rotas e domic\'ilios; garantir supervis\~ao no campo; sincronizar e verificar dados diariamente. \\
\midrule
\textbf{CATI} &
Custo moderado; rapidez e amplo alcance geogr\'afico; maior anonimato em temas sens\'iveis; f\'acil supervis\~ao da equipe. &
Exclui popula\c{c}\~oes sem telefone; question\'arios devem ser curtos e objetivos; taxas de recusa mais elevadas. &
Construir base representativa de n\'umeros; adaptar question\'ario para formato telef\^onico; treinar equipe em comunica\c{c}\~ao por voz; definir protocolos de liga\c{c}\~ao; monitorar entrevistas em tempo real. \\
\midrule
\textbf{Online} &
Custo muito baixo; rapidez e alcance global; conveni\^encia para a pessoa respondente; permite l\'ogica automatizada (\textit{skips}, randomiza\c{c}\~ao). &
Vi\'es de cobertura (exclus\~ao digital); baixas taxas de resposta; pouco controle sobre quem responde. &
Selecionar plataforma online adequada; otimizar \textit{design} para diferentes dispositivos; realizar pr\'e-testes completos; implementar estrat\'egias de divulga\c{c}\~ao e recrutamento; enviar lembretes e monitorar respostas. \\
\bottomrule
\end{tabular}
\figmeta{Elabora\c{c}\~ao pr\'opria.}{}{}
\end{table}

A escolha do m\'etodo quantitativo envolve, em qualquer cen\'ario, um equil\'ibrio entre rigor cient\'ifico, viabilidade operacional e \'etica. Pr\'e-teste dos instrumentos, treinamento das equipes e monitoramento da qualidade s\~ao denominadores comuns: valem para qualquer modalidade \citep{morra2009road,kopper2021survey}.

\section{Metodologias de coleta qualitativa}

Nas avalia\c{c}\~oes de pol\'iticas p\'ublicas, dados qualitativos d\~ao conta de coisas que dados num\'ericos t\^ipicamente n\~ao d\~ao: examinar intera\c{c}\~oes entre atores estatais e n\~ao estatais, identificar padr\~oes de tomada de decis\~ao, analisar a tens\~ao entre regras formais e normas socioculturais. S\~ao eles que revelam o \textit{como} e o \textit{porqu\^e} dos fen\^omenos --- a din\^amica de implementa\c{c}\~ao, as barreiras pr\'aticas que os n\'umeros n\~ao captam, a experi\^encia concreta de quem participa ou implementa. Permitem ainda detectar efeitos n\~ao previstos, distinguir falhas de teoria de falhas de implementa\c{c}\~ao e capturar mudan\c{c}as estruturais que indicadores convencionais dificilmente mensuram \citep{fgvclear2025mea}.

As entrevistas s\~ao o instrumento qualitativo mais usado. Permitem captar nuances do processo de implementa\c{c}\~ao, identificar desafios e analisar como as diretrizes institucionais se traduzem nas din\^amicas cotidianas \citep{lotta2020categorizando}. Em contraste com os \textit{surveys}, t\^em estrutura\c{c}\~ao menor e abrem espa\c{c}o \`a espontaneidade das respostas, o que favorece a emerg\^encia de temas e interpreta\c{c}\~oes n\~ao antecipados no desenho da pesquisa \citep{lima2016entrevistas}.

A entrevista \'e uma intera\c{c}\~ao comunicativa: a pessoa pesquisadora conduz o di\'alogo a partir de um roteiro (mais ou menos estruturado, conforme os objetivos do estudo) e a pessoa entrevistada contribui com relatos, percep\c{c}\~oes e interpreta\c{c}\~oes sobre sua experi\^encia. \'E a comunica\c{c}\~ao direta que viabiliza explorar as formas de agir, pensar e significar dos sujeitos em contextos sociais e institucionais espec\'ificos \citep{leitao2024entrevista}.

Na coleta de dados qualitativos, a escolha do grau de estrutura\c{c}\~ao das entrevistas depende de fatores como o n\'ivel de conhecimento pr\'evio sobre o tema investigado, o n\'umero e o perfil das pessoas participantes, bem como a articula\c{c}\~ao entre a fundamenta\c{c}\~ao te\'orica e a produ\c{c}\~ao de evid\^encias emp\'iricas. Quanto maior o conhecimento acumulado e a necessidade de comparabilidade entre respostas, mais estruturada tende a ser a ferramenta. Por outro lado, quando se busca explorar percep\c{c}\~oes, significados ou fen\^omenos pouco conhecidos, privilegia-se maior flexibilidade nos instrumentos.

Uma modalidade particularmente relevante nesse sentido \'e a entrevista com informantes-chave (KII, do ingl\^es \textit{Key Informant Interview}), na qual as pessoas interlocutoras s\~ao selecionadas intencionalmente por deterem conhecimento especializado, posi\c{c}\~ao privilegiada ou experi\^encia direta sobre o fen\^omeno investigado. Nesse tipo de entrevista, a pessoa entrevistada \'e tratada como especialista, e o di\'alogo visa acessar n\~ao apenas experi\^encias individuais, mas tamb\'em perspectivas institucionais, din\^amicas organizacionais e fatores contextuais que dificilmente seriam captados por instrumentos quantitativos \citep{bryman2012social,neuman2014social}.

As t\'ecnicas de entrevista podem ser presenciais ou virtuais. As presenciais favorecem v\'inculo e confian\c{c}a, al\'em de dar aten\c{c}\~ao aos aspectos n\~ao verbais da comunica\c{c}\~ao. As virtuais ampliam o alcance da pesquisa --- viabilizam participa\c{c}\~ao geograficamente dispersa, reduzem custos e tempo de deslocamento --- mas trazem suas pr\'oprias limita\c{c}\~oes: instabilidades de conex\~ao, privacidade do ambiente da pessoa entrevistada, menor capta\c{c}\~ao de express\~oes n\~ao verbais.

\subsection{Entrevistas estruturadas}

Na entrevista estruturada, a pessoa pesquisadora segue um roteiro previamente definido, com perguntas fixas e enunciado padronizado, apresentadas a todas as entrevistadas na mesma ordem. O formato \'e r\'igido por desenho: \citet{silva2019aplicacao} observam que a padroniza\c{c}\~ao limita a flexibilidade da pesquisadora durante a intera\c{c}\~ao.

O objetivo \'e reduzir a influ\^encia subjetiva da entrevistadora, aumentar a comparabilidade entre participantes e tornar a coleta replic\'avel --- pontos especialmente relevantes em pesquisas de maior escala ou conduzidas por equipes amplas. O custo dessa padroniza\c{c}\~ao \'e a profundidade explorat\'oria: as respostas seguem op\c{c}\~oes preestabelecidas ou um script fixo, aproximando a t\'ecnica de um question\'ario aplicado oralmente.

As preocupa\c{c}\~oes metodol\'ogicas centrais s\~ao tr\^es: coer\^encia entre roteiro, hip\'oteses e objetivos; defini\c{c}\~ao adequada das perguntas e das op\c{c}\~oes de resposta; e neutralidade da pessoa pesquisadora durante a coleta. Mesmo no formato estruturado h\'a alguma margem para que ela adapte vocabul\'ario e linguagem ao contexto da intera\c{c}\~ao, favorecendo a constru\c{c}\~ao de sentido compartilhado \citep{silva2006entrevista}. Respostas registradas em formato textual costumam ser depois organizadas, categorizadas e submetidas a an\'alise de conte\'udo ou a procedimentos quantitativos, em busca de padr\~oes e recorr\^encias nas falas \citep{fraser2004fala}.

Segundo \citet{patton1990qualitative}, a fun\c{c}\~ao da entrevista estruturada \'e assegurar que todas as participantes recebam est\'imulos id\^enticos --- o que aumenta confiabilidade e comparabilidade, com o trade-off j\'a mencionado.

\paragraph{Vantagens.} A padroniza\c{c}\~ao garante que todas as pessoas entrevistadas recebam as mesmas perguntas na mesma ordem e forma, reduzindo vieses de aplica\c{c}\~ao e gerando alta comparabilidade entre respostas. A uniformidade aumenta a confiabilidade dos dados, viabiliza an\'alises estat\'isticas (sobretudo quando h\'a perguntas fechadas) e torna a coleta mais r\'apida e previs\'ivel, o que pesa em pesquisas de grande escala. A formula\c{c}\~ao uniforme reduz tamb\'em as varia\c{c}\~oes de interpreta\c{c}\~ao entre entrevistadas.

\paragraph{Desvantagens.} O custo da rigidez aparece em v\'arias frentes. A pessoa entrevistadora n\~ao pode adaptar perguntas nem aprofundar respostas relevantes que surjam durante a entrevista, o que limita a capta\c{c}\~ao de nuances e significados subjetivos. Qualquer ambiguidade na reda\c{c}\~ao se reproduz em todas as entrevistas, comprometendo a validade. Perguntas fechadas ou excessivamente direcionadas tendem a induzir respostas padronizadas e pouco informativas; e o formato r\'igido n\~ao favorece a compreens\~ao de fatores situacionais ou explica\c{c}\~oes mais amplas. No conjunto, a t\'ecnica privilegia comparabilidade e mensura\c{c}\~ao em detrimento das dimens\~oes interpretativas e simb\'olicas do fen\^omeno.

\subsection{Entrevistas semiestruturadas}

A entrevista semiestruturada parte de um roteiro pr\'e-elaborado, mas confere \`a pessoa pesquisadora margem para ajustar ordem, formula\c{c}\~ao e profundidade das perguntas conforme o di\'alogo se desenvolve. \'E o ponto de equil\'ibrio entre a padroniza\c{c}\~ao necess\'aria \`a comparabilidade e a abertura que captura a riqueza das narrativas. Da\'i ser uma das t\'ecnicas qualitativas mais usadas em pesquisas que buscam compreender significados, percep\c{c}\~oes, pr\'aticas e experi\^encias.

A elabora\c{c}\~ao do roteiro come\c{c}a pela defini\c{c}\~ao de temas, eixos anal\'iticos ou eventos significativos relacionados ao objeto de estudo. Esses elementos orientam a reconstru\c{c}\~ao de experi\^encias, trajet\'orias e percep\c{c}\~oes das pessoas interlocutoras \citep{alonso2016metodos}. O roteiro se organiza em torno de quest\~oes principais, que funcionam como eixos da intera\c{c}\~ao, mas habilitam a entrevistadora a aprofundar aspectos emergentes ou introduzir novas perguntas conforme a din\^amica avan\c{c}a \citep{lima2016entrevistas}. Essa abertura \'e o que viabiliza captar dimens\~oes n\~ao previstas no planejamento inicial.

A pesquisadora deve garantir que todos os t\'opicos previamente definidos sejam abordados, mas a sequ\^encia, a reda\c{c}\~ao e a \^enfase das perguntas podem se ajustar ao contexto de cada entrevista. Essa maleabilidade \'e o que torna o m\'etodo adequado \`a investiga\c{c}\~ao de significados, motiva\c{c}\~oes e interpreta\c{c}\~oes --- preservando, ao mesmo tempo, a coer\^encia com os objetivos, hip\'oteses e categorias anal\'iticas da pesquisa.

Uma estrat\'egia recorrente \'e conduzir a entrevista de forma progressiva: perguntas mais abertas e explorat\'orias no in\'icio, voltadas \`a livre express\~ao da participante; perguntas mais espec\'ificas ou direcionadas \`a medida que a conversa avan\c{c}a \citep{silva2006entrevista}. O movimento ajuda a estabelecer confian\c{c}a, estimula a narrativa do entrevistado e, simultaneamente, assegura a obten\c{c}\~ao de informa\c{c}\~oes relevantes para os objetivos do estudo.

\paragraph{Vantagens.} A flexibilidade na condu\c{c}\~ao permite \`a entrevistadora adaptar sequ\^encia, aprofundamento e formula\c{c}\~ao das perguntas conforme as respostas, favorecendo intera\c{c}\~ao mais natural. O formato combina comparabilidade m\'inima entre entrevistas com liberdade para explorar temas emergentes, sem perder foco anal\'itico. Permite compreender percep\c{c}\~oes, valores, motiva\c{c}\~oes e experi\^encias com profundidade e contextualiza\c{c}\~ao, e a abertura metodol\'ogica favorece a identifica\c{c}\~ao de aspectos n\~ao previstos no roteiro. \'E adapt\'avel a diferentes contextos e perfis de participantes, do gestor s\^enior \`a benefici\'aria de programa.

\paragraph{Desvantagens.} A flexibilidade dificulta a compara\c{c}\~ao direta entre entrevistas e, sem controle metodol\'ogico cuidadoso, reduz a consist\^encia anal\'itica. A qualidade da entrevista depende fortemente da entrevistadora --- exige preparo t\'ecnico, empatia e dom\'inio do tema para conduzir sem enviesar. A presen\c{c}a da pesquisadora pode influenciar respostas, e a interpreta\c{c}\~ao depende da sua sensibilidade anal\'itica. O car\'ater interativo dificulta a replicabilidade entre contextos. E roteiros mal estruturados geram lacunas tem\'aticas ou dispers\~ao, comprometendo a coer\^encia da coleta.

\subsection{Entrevistas n\~ao estruturadas}

A entrevista n\~ao estruturada \'e flex\'ivel por desenho. N\~ao h\'a roteiro fixo: os temas emergem durante a intera\c{c}\~ao, e a pesquisadora acompanha o fluxo da conversa, aprofundando pontos de interesse \`a medida que aparecem. O objetivo \'e maximizar o car\'ater explorat\'orio da coleta, favorecendo a identifica\c{c}\~ao de significados, percep\c{c}\~oes e dimens\~oes do fen\^omeno que dificilmente seriam antecipados.

Mesmo sem roteiro estruturado, a pesquisadora se orienta por temas ou eixos gerais que devem ser cobertos ao longo da conversa --- eles funcionam como refer\^encia para formular perguntas durante a entrevista. Tanto a entrevistadora quanto a pessoa entrevistada t\^em ampla liberdade, e \'e comum que a participante introduza t\'opicos n\~ao previstos no desenho original \citep{lima2016entrevistas}.

A abordagem cai bem em estudos explorat\'orios, em que h\'a pouco ou nenhum conhecimento acumulado sobre o tema. Nesses casos, a entrevistadora formula perguntas amplas e abertas, que estimulam as participantes a discorrer livremente \citep{silva2019aplicacao}. A postura predominante \'e de ``abertura'', com interfer\^encias pontuais para que a entrevistada aprofunde um tema ou descreva um acontecimento espec\'ifico \citep{fraser2004fala}.

As perguntas emergem do contexto interacional, n\~ao s\~ao predeterminadas \citep{jimenez2012entrevista}. O foco recai sobre a explora\c{c}\~ao detalhada das percep\c{c}\~oes, experi\^encias e significados atribu\'idos pelas participantes. Isso exige escuta atenta, acompanhamento da narrativa e aprofundamento sucessivo --- a entrevistadora clarifica ambiguidades, explora nuances do discurso e procura entender como as interlocutoras constroem sentido sobre o fen\^omeno investigado. O resultado \'e uma an\'alise densa e contextualizada.

\paragraph{Vantagens.} M\'axima flexibilidade metodol\'ogica: a entrevistadora pode adaptar totalmente a condu\c{c}\~ao da conversa, explorando temas que emergem espontaneamente. O formato propicia compreens\~ao profunda e contextualizada das percep\c{c}\~oes, experi\^encias e significados atribu\'idos pelas participantes, e \'e a t\'ecnica que mais favorece a emerg\^encia de dimens\~oes anal\'iticas n\~ao antecipadas. O ambiente aberto encoraja respostas espont\^aneas e reflexivas. \'E especialmente \'util em estudos explorat\'orios sem hip\'oteses consolidadas, e a aus\^encia de roteiro fixo permite ajuste fluido a diferentes perfis de entrevistadas.

\paragraph{Desvantagens.} A aus\^encia de padroniza\c{c}\~ao dificulta a compara\c{c}\~ao sistem\'atica entre entrevistas e a constru\c{c}\~ao de categorias comuns de an\'alise. A condu\c{c}\~ao depende intensamente da experi\^encia da pesquisadora --- manter foco tem\'atico e explorar respostas sem direcion\'a-las exige treino. A natureza aberta e heterog\^enea das respostas exige procedimentos de codifica\c{c}\~ao e interpreta\c{c}\~ao mais complexos e demorados. A singularidade de cada intera\c{c}\~ao torna a replicabilidade praticamente imposs\'ivel. A entrevistadora pode projetar suas pr\'oprias percep\c{c}\~oes sobre o discurso da entrevistada, agravando vieses interpretativos. Por fim, transcri\c{c}\~ao e an\'alise s\~ao trabalhosas, o que costuma limitar o n\'umero de participantes.

\subsection{Entrevistas em grupo e grupos focais}

Entrevistas em grupo e grupos focais s\~ao frequentemente tratados como sin\^onimos, mas n\~ao s\~ao. Embora ambos envolvam coleta com m\'ultiplas participantes simultaneamente, distinguem-se em tr\^es aspectos: especificidade tem\'atica, interesse pela intera\c{c}\~ao em si e papel atribu\'ido \`as participantes na produ\c{c}\~ao de conhecimento \citep{bryman2012social,neuman2014social}.

A entrevista em grupo aplica perguntas a um conjunto de participantes reunidas ao mesmo tempo, em busca de respostas individuais simult\^aneas. O escopo tem\'atico pode ser amplo, e a t\'ecnica costuma ser adotada por economia de tempo e recursos --- a pesquisadora est\'a mais interessada na efici\^encia da coleta do que nas din\^amicas relacionais do grupo \citep{bryman2012social}.

O grupo focal \'e mais especializado. Caracteriza-se pela delimita\c{c}\~ao de um tema espec\'ifico e pelo interesse expl\'icito na intera\c{c}\~ao entre participantes como fonte central de dados. A intera\c{c}\~ao n\~ao \'e efeito colateral da reuni\~ao em grupo --- ela \'e o pr\'oprio objeto de aten\c{c}\~ao anal\'itica: o di\'alogo permite n\~ao apenas a express\~ao de percep\c{c}\~oes individuais, mas tamb\'em a negocia\c{c}\~ao de significados, a manifesta\c{c}\~ao de consensos e a explicita\c{c}\~ao de dissensos. \'E a t\'ecnica adequada para captar din\^amicas sociais, explorar representa\c{c}\~oes compartilhadas e compreender como opini\~oes se constroem e se transformam no contexto da intera\c{c}\~ao grupal.

A discuss\~ao \'e orientada por uma pergunta principal e por perguntas secund\'arias \citep{almeida2016grupos}. As participantes s\~ao incentivadas a argumentar, questionar e desafiar perspectivas umas das outras, o que pode lev\'a-las a revisar suas pr\'oprias opini\~oes ao longo da discuss\~ao. Esse processo de constru\c{c}\~ao coletiva de significado tende a produzir relatos mais realistas do pensamento do grupo do que entrevistas individuais. A pesquisadora atua como mediadora --- planeja o roteiro, define perguntas disparadoras e realiza interven\c{c}\~oes pontuais para manter participa\c{c}\~ao equilibrada, est\'imulo \`a intera\c{c}\~ao e foco na tem\'atica \citep{jimenez2012entrevista,bryman2012social}.

A vantagem distintiva do grupo focal em rela\c{c}\~ao \`as entrevistas individuais est\'a justamente em captar essas din\^amicas \citep{silva2021modera}. Para que a t\'ecnica funcione, os grupos devem ser pequenos --- geralmente entre 6 e 12 participantes ---, o que viabiliza participa\c{c}\~ao ativa de todas, acompanhamento detalhado das intera\c{c}\~oes e monitoramento de assimetrias de poder (l\'ideres de opini\~ao dominando o debate ou inibindo pessoas mais reticentes).

Para al\'em de evitar pessoas com rela\c{c}\~ao hier\'arquica direta entre si, o \textit{status} importa de forma mais ampla: misturar n\'iveis socioecon\^omicos ou posi\c{c}\~oes de autoridade muito distintos pode levar as participantes a expressar apenas vis\~oes culturalmente esperadas ou a se conformar \`a opini\~ao da maioria por press\~ao social. A composi\c{c}\~ao do grupo \'e decis\~ao deliberada: heterogeneidade quando se quer captar variedade de perspectivas, homogeneidade quando se quer explorar um tema espec\'ifico em profundidade dentro de um grupo com experi\^encias semelhantes. Em contrapartida, o formato coletivo pode empoderar participantes --- diminui a dist\^ancia entre pesquisadora e pesquisadas, transfere parte do controle da conversa para o grupo e permite que vozes marginalizadas encontrem suporte m\'utuo.

A an\'alise dos dados de grupos focais n\~ao se reduz a contar respostas. Converg\^encia e diverg\^encia de opini\~oes, \^enfase emocional, linguagem corporal --- todos esses elementos s\~ao informa\c{c}\~ao sobre as percep\c{c}\~oes coletivas e individuais e precisam entrar na interpreta\c{c}\~ao.

\paragraph{Vantagens.} Permite observar intera\c{c}\~oes, negocia\c{c}\~oes de significados, consensos e dissensos --- t\'ipico do que entrevistas individuais n\~ao capturam. Facilita a identifica\c{c}\~ao de percep\c{c}\~oes compartilhadas e o entendimento de como opini\~oes se constroem no contexto grupal. A intera\c{c}\~ao entre participantes encoraja a manifesta\c{c}\~ao de ideias que n\~ao surgiriam num formato individual. A mediadora pode ajustar interven\c{c}\~oes e aprofundar t\'opicos emergentes. Em termos pr\'aticos, m\'ultiplas perspectivas s\~ao obtidas em uma \'unica sess\~ao, com ganho de tempo em rela\c{c}\~ao a entrevistas individuais. E o formato permite observar comportamentos n\~ao verbais --- linguagem corporal, entona\c{c}\~ao, intera\c{c}\~oes --- que enriquecem a interpreta\c{c}\~ao.

\paragraph{Desvantagens.} A influ\^encia social \'e o ponto mais delicado: participantes podem se sentir pressionadas a concordar com opini\~oes dominantes ou a se autocensurar. O sucesso depende fortemente da habilidade da mediadora em equilibrar est\'imulo \`a participa\c{c}\~ao com manuten\c{c}\~ao do foco. Indiv\'iduos mais assertivos tendem a monopolizar a conversa, enquanto outros recuam. O planejamento log\'istico \'e complexo --- coordenar hor\'ario, local, recrutamento e composi\c{c}\~ao do grupo n\~ao \'e trivial. A interpreta\c{c}\~ao precisa lidar simultaneamente com conte\'udo, intera\c{c}\~ao e contexto, o que aumenta a complexidade anal\'itica. E h\'a o problema do tamanho: grupos pequenos reduzem diversidade; grupos grandes dificultam participa\c{c}\~ao ativa e registro detalhado.

\section{Metodologias mistas de coleta}

A metodologia mista n\~ao consiste simplesmente em coletar dados de dois tipos em paralelo. Ela pressup\~oe uma decis\~ao metodol\'ogica deliberada sobre como os componentes se articulam --- qual deles prevalece, em que momento se realiza cada coleta e como os dados s\~ao integrados na an\'alise \citep{creswell2014research}. Quando planejada com rigor, a abordagem mista permite que os resultados quantitativos orientem a sele\c{c}\~ao de casos para o aprofundamento qualitativo, ou que os achados qualitativos sinalizem hip\'oteses a serem testadas em \textit{surveys} de maior escala.

\subsection{Principais arranjos de coleta mista}

Tr\^es arranjos s\~ao particularmente frequentes em avalia\c{c}\~oes de pol\'iticas p\'ublicas:

\paragraph{Sequencial explorat\'orio.} A coleta qualitativa antecede a quantitativa. Entrevistas ou grupos focais explorat\'orios informam a constru\c{c}\~ao do question\'ario --- identificando dimens\~oes relevantes, terminologia adequada ao p\'ublico-alvo e categorias de resposta. \'E especialmente \'util quando o fen\^omeno \'e pouco conhecido ou quando a linguagem dos instrumentos precisa ser ajustada a contextos espec\'ificos.

\paragraph{Sequencial explanat\'orio.} O \textit{survey} \'e realizado primeiro; os achados quantitativos orientam a sele\c{c}\~ao de casos para entrevistas ou grupos focais de aprofundamento. Esse arranjo \'e comum em avalia\c{c}\~oes de impacto que identificam subgrupos com resultados inesperados e buscam compreender as raz\~oes por tr\'as dessas diferen\c{c}as. A vantagem \'e que os dados qualitativos s\~ao coletados com um objetivo preciso --- interpretar o que os n\'umeros n\~ao explicam.

\paragraph{Concomitante (triangula\c{c}\~ao).} \textit{Surveys} e entrevistas s\~ao aplicados simultaneamente, com integra\c{c}\~ao na fase de an\'alise. A triangula\c{c}\~ao permite verificar a converg\^encia entre fontes, identificar contradi\c{c}\~oes e obter uma descri\c{c}\~ao mais completa do fen\^omeno \citep{patton2002qualitative,creswell2014research}. Exige maior coordena\c{c}\~ao de campo e maior esfor\c{c}o anal\'itico, mas \'e o arranjo que mais potencializa a complementaridade entre os m\'etodos.

\subsection{Crit\'erios para a escolha e o sequenciamento dos instrumentos}

A op\c{c}\~ao pela metodologia mista deve ser guiada pelas perguntas de pesquisa e pelos recursos dispon\'iveis --- e n\~ao pela ideia de que ``quanto mais m\'etodos, melhor''. As seguintes quest\~oes pr\'aticas orientam a decis\~ao:

\begin{itemize}
\item \textbf{Qual \'e a pergunta que o m\'etodo adicional responde?} Cada instrumento deve ter uma fun\c{c}\~ao anal\'itica clara --- explorar, confirmar, aprofundar ou triangular. Instrumentos acrescentados sem fun\c{c}\~ao definida aumentam o custo e a complexidade da coleta sem necessariamente agregar valor aos resultados.
\item \textbf{Quando e em que ordem coletar?} A sequ\^encia afeta tanto a log\'istica de campo quanto a qualidade das informa\c{c}\~oes. A coleta qualitativa pr\'evia reduz erros de mensura\c{c}\~ao no \textit{survey}; a qualitativa posterior permite explorar resultados inesperados. A coleta simult\^anea exige coordena\c{c}\~ao rigorosa entre as equipes e protocolos comuns de controle de qualidade.
\item \textbf{Como integrar os dados na an\'alise?} A integra\c{c}\~ao pode ocorrer na interpreta\c{c}\~ao (comparar resultados separadamente), na transforma\c{c}\~ao (converter dados qualitativos em c\'odigos quantific\'aveis) ou na conex\~ao (usar os resultados de um m\'etodo para orientar a an\'alise do outro). Cada estrat\'egia tem implica\c{c}\~oes para o relato dos resultados e deve ser descrita com transpar\^encia metodol\'ogica \citep{creswell2014research,bryman2012social}.
\end{itemize}

\subsection{Desafios operacionais e boas pr\'aticas}

A combina\c{c}\~ao de instrumentos traz desafios que v\~ao al\'em do desenho metodol\'ogico. Na pr\'atica, equipes com perfis diferentes --- pessoas pesquisadoras quantitativas e entrevistadoras qualitativas --- precisam coordenar protocolos, cronogramas e padr\~oes de registro. Alguns cuidados s\~ao especialmente relevantes:

\begin{itemize}
\item \textbf{Evitar a sobrecarga de pessoas participantes:} aplicar \textit{survey} e entrevista \`as mesmas pessoas exige cuidado com o tempo demandado e com o intervalo entre as coletas. Quando poss\'ivel, recomenda-se que a coleta qualitativa seja realizada com uma subamostra intencionalmente selecionada, e n\~ao com todas as pessoas participantes do \textit{survey}.
\item \textbf{Registrar as decis\~oes de integra\c{c}\~ao:} como os dados de diferentes instrumentos foram combinados deve constar no relat\'orio metodol\'ogico. A aus\^encia de transpar\^encia sobre esse ponto \'e uma das cr\'iticas mais comuns \`as pesquisas de m\'etodos mistos \citep{bryman2012social}.
\item \textbf{Planejar o or\c{c}amento de forma integrada:} coletas mistas t\^em custo significativamente maior do que coletas com instrumento \'unico. O planejamento deve prever os recursos espec\'ificos de cada componente, incluindo transcri\c{c}\~ao e an\'alise dos dados qualitativos, frequentemente subestimados na fase de elabora\c{c}\~ao do projeto \citep{duru2021grant}.
\end{itemize}

Bem executada, a metodologia mista combina o que cada m\'etodo isolado n\~ao entrega sozinho --- os dados quantitativos respondem ``quanto''; os qualitativos respondem ``por qu\^e''.

% ============================================================
% 3. AMOSTRA E SELEÇÃO DE PESSOAS PARTICIPANTES
% ============================================================
\chapter{Amostra e sele\c{c}\~ao de pessoas participantes}

A sele\c{c}\~ao das participantes determina, de sa\'ida, a qualidade, a validade e a representatividade da pesquisa. Em estudos quantitativos e qualitativos, o desenho da amostra define a credibilidade dos achados e o que se pode generalizar a partir deles \citep{neuman2014social,gertler2018impact}.\footnote{A generaliza\c{c}\~ao em estudos qualitativos \'e tema de debate na literatura. \citet{mattos2011resultados} sistematiza a discuss\~ao.}

Mas a amostragem n\~ao \'e s\'o uma decis\~ao t\'ecnica. \'E tamb\'em \'etica e institucional. O modo como participantes s\~ao inclu\'idas ou exclu\'idas tem implica\c{c}\~oes diretas sobre a equidade e a legitimidade das conclus\~oes --- da\'i o \textit{Guia CLEAR de Avalia\c{c}\~ao de Impacto} \citep{fgvclear2025ai} e o \textit{Guia CLEAR de Monitoramento e Avalia\c{c}\~ao} \citep{fgvclear2025mea} tratarem a amostragem como instrumento de governan\c{c}a, articulando decis\~oes metodol\'ogicas \`as restri\c{c}\~oes operacionais, de tempo e or\c{c}amento.

Em paralelo, tamanho e estrutura da amostra determinam a capacidade do estudo de detectar efeitos reais. Amostras subdimensionadas falham em identificar impactos relevantes; superdimensionadas desperdi\c{c}am recursos. O \textit{An\'alise de Poder Estat\'istico} \citep{fgvclear2025poder} aprofunda esse trade-off --- o planejamento amostral precisa equilibrar precis\~ao estat\'istica, viabilidade operacional e efici\^encia, e essa decis\~ao entra no desenho da pesquisa desde o in\'icio.

\begin{conexaobox}
Este cap\'itulo dialoga diretamente com o \textit{An\'alise de Poder Estat\'istico}, que detalha o c\'alculo de tamanho amostral m\'inimo, MDE (\textit{Minimum Detectable Effect}) e \textit{design effect}, e com o \textit{Guia CLEAR de Avalia\c{c}\~ao de Impacto}, que discute aleatoriza\c{c}\~ao e desenhos experimentais.
\end{conexaobox}

A literatura distingue dois grandes grupos de m\'etodos: \textbf{amostragem probabil\'istica}, voltada \`a infer\^encia estat\'istica e generaliza\c{c}\~ao; e \textbf{amostragem n\~ao probabil\'istica}, mais adequada \`a explora\c{c}\~ao qualitativa e compreens\~ao de contextos \citep{bryman2012social,neuman2014social}.

\section{M\'etodos comuns de amostragem probabil\'istica}

Na amostragem probabil\'istica, cada elemento da popula\c{c}\~ao tem probabilidade conhecida e diferente de zero de ser selecionado --- normalmente por meio de um sorteio. Essa propriedade \'e o que viabiliza a infer\^encia estat\'istica: os resultados da amostra podem ser generalizados para a popula\c{c}\~ao, e o erro amostral (diferen\c{c}a esperada entre estimativa amostral e valor real na popula\c{c}\~ao) pode ser quantificado. \'E o padr\~ao-ouro para estudos quantitativos e avalia\c{c}\~oes experimentais \citep{gertler2018impact}.

\subsubsection*{1) Amostragem aleat\'oria simples (\textit{simple random sampling})}

\'E a forma mais b\'asica de amostragem probabil\'istica, onde cada unidade (ex.: indiv\'iduo, domic\'ilio) na lista da popula\c{c}\~ao (\textit{frame}) tem exatamente a mesma chance de ser selecionada. Normalmente, utiliza-se uma tabela de n\'umeros aleat\'orios ou um gerador de n\'umeros aleat\'orios para realizar a sele\c{c}\~ao.

\begin{itemize}
\item \textbf{Vantagem:} \'e conceitualmente simples e f\'acil de implementar para popula\c{c}\~oes pequenas e homog\^eneas, pois requer apenas uma lista completa de todos os elementos da popula\c{c}\~ao.
\item \textbf{Desvantagem:} requer lista completa da popula\c{c}\~ao (\textit{sampling frame}), tornando-se logisticamente invi\'avel em contextos amplos. Em popula\c{c}\~oes grandes e dispersas, a necessidade de identificar, listar e acessar todos os elementos tende a aumentar significativamente a complexidade operacional e os custos do processo \citep{bryman2012social,neuman2014social}.
\item \textbf{Exemplo:} sorteio de domic\'ilios numerados em um \'unico munic\'ipio.
\end{itemize}

\subsubsection*{2) Amostragem estratificada (\textit{stratified sampling})}

Nesta abordagem, a popula\c{c}\~ao \'e, primeiro, dividida em subgrupos homog\^eneos e mutuamente exclusivos, chamados de estratos (ex.: por regi\~ao geogr\'afica, tamanho da propriedade, faixa et\'aria, g\^enero). Em seguida, uma amostra aleat\'oria simples \'e extra\'ida de cada estrato independentemente.

\begin{itemize}
\item \textbf{Vantagem:} garante a representatividade de subgrupos-chave na amostra final, o que \'e crucial para an\'alises comparativas. Al\'em disso, frequentemente aumenta a precis\~ao estat\'istica (reduz a vari\^ancia) das estimativas, pois a variabilidade dentro de cada estrato \'e menor do que na popula\c{c}\~ao como um todo \citep{neuman2014social,iarossi2006power,gertler2018impact}.
\item \textbf{Desvantagem:} exige conhecimento pr\'evio da popula\c{c}\~ao para definir os estratos relevantes e obter informa\c{c}\~oes para classificar cada elemento.
\item \textbf{Exemplo:} sorteio de escolas separadamente por zona urbana/rural e porte.
\end{itemize}

\subsubsection*{3) Amostragem por conglomerados (\textit{cluster sampling})}

Aqui, reconhece-se que a popula\c{c}\~ao \'e naturalmente dividida em grupos, ou conglomerados (ex.: vilas, escolas, quarteir\~oes). Com isso em mente, primeiro, uma amostra aleat\'oria desses conglomerados \'e selecionada. Em seguida, todos os elementos dentro dos conglomerados escolhidos s\~ao pesquisados (amostragem em um est\'agio) ou uma amostra aleat\'oria de elementos dentro de cada conglomerado \'e extra\'ida (amostragem em dois est\'agios).

\begin{itemize}
\item \textbf{Vantagem:} \'e altamente eficiente em termos de custo e log\'istica, especialmente quando a popula\c{c}\~ao est\'a geograficamente dispersa. Visitar alguns vilarejos e entrevistar todas as pessoas residentes tende a ser menos custoso do que visitar domic\'ilios espalhados por uma regi\~ao vasta \citep{iarossi2006power,neuman2014social}.
\item \textbf{Desvantagem:} para um mesmo tamanho de amostra, tende a ser menos preciso do que a amostragem aleat\'oria simples ou estratificada, pois elementos dentro de um conglomerado tendem a ser similares entre si (homogeneidade intra-conglomerado). Tamb\'em requer ajuste para efeito de desenho (\textit{design effect}) \citep{gertler2018impact}.
\item \textbf{Exemplo:} sele\c{c}\~ao de munic\'ipios e, dentro de cada um, sorteio de escolas p\'ublicas.
\end{itemize}

\subsubsection*{4) Amostragem com probabilidade proporcional ao tamanho (PPS)}

Trata-se de uma variante sofisticada da amostragem por conglomerados, usada quando os conglomerados diferem de tamanho. Cada conglomerado tem probabilidade proporcional ao n\'umero de unidades que cont\'em. Nela, conglomerados maiores (ex.: vilas com mais habitantes) t\^em uma probabilidade maior de serem selecionados na primeira etapa do que conglomerados menores. Este m\'etodo \'e combinado com a sele\c{c}\~ao de um n\'umero fixo de elementos dentro de cada conglomerado selecionado.

\begin{itemize}
\item \textbf{Vantagem:} \'e essencial para garantir que cada indiv\'iduo na popula\c{c}\~ao tenha uma chance igual de ser selecionado, mesmo quando os conglomerados t\^em tamanhos muito diferentes. Sem o PPS, indiv\'iduos em conglomerados grandes teriam uma chance menor de serem selecionados do que aqueles em conglomerados pequenos \citep{iarossi2006power,neuman2014social}.
\item \textbf{Desvantagem:} essa amostragem \'e mais complexa de implementar, pois requer o conhecimento do tamanho de cada conglomerado na popula\c{c}\~ao para calcular suas probabilidades de sele\c{c}\~ao.
\item \textbf{Exemplo:} em uma pesquisa nacional, munic\'ipios com mais alunas e alunos t\^em maior probabilidade de inclus\~ao \citep{iarossi2006power,neuman2014social}.
\end{itemize}

A Tabela~\ref{tab:tipos_amostragem} sintetiza os tipos de amostragem mais comuns, suas aplica\c{c}\~oes, vantagens e limita\c{c}\~oes, servindo como refer\^encia pr\'atica para a etapa de planejamento da coleta.

\begin{table}[H]
\centering
\caption{Tipos de amostragem e principais caracter\'isticas}
\label{tab:tipos_amostragem}
\small
\begin{tabular}{p{0.15\linewidth} p{0.20\linewidth} p{0.18\linewidth} p{0.18\linewidth} p{0.20\linewidth}}
\toprule
\textbf{Tipo} & \textbf{Descri\c{c}\~ao} & \textbf{Vantagens} & \textbf{Limita\c{c}\~oes} & \textbf{Quando utilizar} \\
\midrule
Aleat\'oria simples & Cada unidade tem a mesma probabilidade de ser selecionada. & Simples, transparente e estatisticamente robusta. & Requer lista completa da popula\c{c}\~ao; pouco vi\'avel em popula\c{c}\~oes dispersas. & Popula\c{c}\~oes pequenas e homog\^eneas com cadastro dispon\'ivel. \\
\midrule
Estratificada & Popula\c{c}\~ao dividida em estratos homog\^eneos e amostras sorteadas em cada um. & Garante representatividade de subgrupos e reduz vari\^ancia. & Exige conhecimento pr\'evio detalhado da popula\c{c}\~ao. & Quando h\'a interesse em comparar subgrupos. \\
\midrule
Por conglomerados & Sele\c{c}\~ao de grupos naturais (vilas, escolas, setores) e posterior sele\c{c}\~ao dentro deles. & Reduz custos e deslocamentos; vi\'avel em grandes \'areas. & Menor precis\~ao; exige ajuste para \textit{design effect}. & Popula\c{c}\~oes amplas e geograficamente dispersas. \\
\midrule
PPS & Conglomerados maiores t\^em maior probabilidade de sele\c{c}\~ao. & Garante probabilidade igual de sele\c{c}\~ao individual. & C\'alculos mais complexos e depend\^encia de informa\c{c}\~oes pr\'evias. & Pesquisas nacionais com unidades de tamanhos distintos. \\
\midrule
Intencional / Conveni\^encia & Sele\c{c}\~ao deliberada ou baseada em acessibilidade. & Simples, \'util em estudos explorat\'orios. & N\~ao permite generaliza\c{c}\~ao estat\'istica. & Pesquisas qualitativas ou contextos com limita\c{c}\~ao de acesso. \\
\bottomrule
\end{tabular}
\figmeta{Elabora\c{c}\~ao pr\'opria.}{}{}
\end{table}

\section{Determina\c{c}\~ao do tamanho da amostra (amostragem probabil\'istica)}

Dimensionar a amostra \'e definir quantas pessoas participar\~ao do estudo. \'E decis\~ao de desenho, n\~ao de execu\c{c}\~ao --- e o c\'alculo malfeito gera dois problemas opostos. Amostras subdimensionadas n\~ao conseguem detectar efeitos que existem na popula\c{c}\~ao. Amostras superdimensionadas s\~ao ineficientes: o custo marginal de cada observa\c{c}\~ao extra n\~ao se justifica pelo ganho infinitesimal de precis\~ao, e o estudo tende a identificar diferen\c{c}as ``estatisticamente significantes'' sem relev\^ancia pr\'atica ou pol\'itica.

\citet{duflo2007using} e \citet{gertler2018impact} discutem o equil\'ibrio necess\'ario entre rigor estat\'istico e viabilidade operacional. Em avalia\c{c}\~oes de pol\'iticas p\'ublicas, isso significa conciliar exig\^encias de precis\~ao com restri\c{c}\~oes log\'isticas, or\c{c}ament\'arias e de tempo de campo.

O \textit{An\'alise de Poder Estat\'istico} \citep{fgvclear2025poder} desenvolve o conceito central a\'i envolvido --- o \textbf{poder estat\'istico}, ou a probabilidade de o estudo detectar um efeito verdadeiro quando ele existe. Estudos com poder menor que 80\% s\~ao considerados pouco confi\'aveis para detec\c{c}\~ao de efeitos reais; estudos com poder muito acima disso operam com amostras maiores do que o necess\'ario.

Cinco par\^ametros principais determinam o tamanho da amostra:

\begin{itemize}
\item \textbf{N\'ivel de confian\c{c}a (geralmente 95\%):} representa o grau de certeza de que o intervalo estimado cont\'em o valor real do par\^ametro populacional. N\'iveis de confian\c{c}a mais altos aumentam a precis\~ao, mas tamb\'em exigem amostras maiores.
\item \textbf{Margem de erro desejada (por exemplo, $\pm$ 5 p.p.):} define o limite de varia\c{c}\~ao aceit\'avel entre o valor observado na amostra e o valor real na popula\c{c}\~ao. Quanto menor a margem de erro, maior dever\'a ser o tamanho da amostra.
\item \textbf{Variabilidade esperada da vari\'avel de interesse:} reflete o quanto os dados s\~ao dispersos. Vari\'aveis com alta variabilidade (por exemplo, renda familiar) exigem amostras maiores do que vari\'aveis mais homog\^eneas (por exemplo, idade em uma popula\c{c}\~ao escolar).
\item \textbf{Efeito de desenho (\textit{design effect}):} ajusta o c\'alculo para o fato de que amostragens complexas (como as por conglomerados) reduzem a precis\~ao estat\'istica em compara\c{c}\~ao \`a amostragem aleat\'oria simples. O efeito de desenho geralmente aumenta o tamanho necess\'ario da amostra em propor\c{c}\~ao \`a homogeneidade intra-conglomerado.
\item \textbf{Taxa de n\~ao resposta:} representa a propor\c{c}\~ao esperada de indiv\'iduos que, por recusa ou aus\^encia, n\~ao participar\~ao da pesquisa. Para compensar essas perdas, recomenda-se inflar o tamanho calculado da amostra inicial.
\end{itemize}

Esses fatores interagem de forma n\~ao linear: um aumento em qualquer um deles (por exemplo, maior n\'ivel de confian\c{c}a ou maior variabilidade esperada) eleva substancialmente o tamanho necess\'ario da amostra. Por isso, o dimensionamento deve ser tratado como um processo iterativo, envolvendo simula\c{c}\~oes e revis\~oes sucessivas.

\textit{Softwares} como Stata, R e Python oferecem fun\c{c}\~oes espec\'ificas para estimar o tamanho da amostra e conduzir an\'alises de poder (\textit{power analysis}). Em avalia\c{c}\~oes experimentais, essa etapa \'e indispens\'avel para determinar o \textbf{tamanho m\'inimo de efeito detect\'avel} (\textit{Minimum Detectable Effect} --- MDE) --- ou seja, o menor impacto que o estudo \'e capaz de identificar com probabilidade estat\'istica aceit\'avel.

O \textit{Guia CLEAR de Avalia\c{c}\~ao de Impacto} \citep{fgvclear2025ai} recomenda que o plano de amostragem contenha documenta\c{c}\~ao sobre os par\^ametros adotados e as decis\~oes tomadas durante esse processo. Especificamente, o documento deve apresentar:

\begin{itemize}
\item Popula\c{c}\~ao e unidade de an\'alise (por exemplo, alunos, escolas, munic\'ipios);
\item Estrat\'egia de sele\c{c}\~ao (probabil\'istica, estratificada, por \textit{cluster}, etc.);
\item Crit\'erios de inclus\~ao e exclus\~ao;
\item Estrutura de estratifica\c{c}\~ao e conglomerados;
\item Justificativa do tamanho amostral, incluindo resultados da an\'alise de poder;
\item Procedimentos para substitui\c{c}\~oes, perdas e controle de n\~ao resposta.
\end{itemize}

Essa documenta\c{c}\~ao formaliza a amostragem como decis\~ao estrat\'egica do desenho avaliativo --- refletindo prioridades de pol\'itica, recursos dispon\'iveis e n\'iveis de precis\~ao aceit\'aveis \citep{fgvclear2025mea,duflo2007using,gertler2018impact}.

\section{Amostragem n\~ao probabil\'istica}

Em estudos qualitativos, a amostragem intencional --- ou n\~ao probabil\'istica --- \'e a estrat\'egia dominante. O objetivo n\~ao \'e representatividade estat\'istica, mas abrang\^encia anal\'itica e profundidade: selecionam-se participantes cujas experi\^encias, posi\c{c}\~oes ou conhecimentos sejam analiticamente significativos em rela\c{c}\~ao ao fen\^omeno investigado \citep{rocha2020teoria}. O que se busca \'e a riqueza e a relev\^ancia da informa\c{c}\~ao produzida. Nas palavras de \citet{bauer2013construcao}, a coleta qualitativa corresponde \`a constru\c{c}\~ao de uma cole\c{c}\~ao finita de materiais textuais e lingu\'isticos que expressam a diversidade das formas de compreens\~ao da realidade social.

A sele\c{c}\~ao intencional envolve etapas sucessivas: identifica\c{c}\~ao e sele\c{c}\~ao das participantes conforme crit\'erios de inclus\~ao e exclus\~ao; convite \`a participa\c{c}\~ao; agendamento das entrevistas \citep{leitao2024entrevista,depaula2014modos}. Em pesquisas de pol\'iticas p\'ublicas, os crit\'erios costumam contemplar diversidade de pap\'eis sociais, trajet\'orias de vida, experi\^encias e posi\c{c}\~oes institucionais. \citet{minayo2017amostragem} sistematiza um conjunto de etapas:

\begin{itemize}
\item Elaborar instrumentos que permitam a compreens\~ao das homogeneidades e diferencia\c{c}\~oes entre grupos;
\item Garantir que a escolha dos grupos estudados contemple as caracter\'isticas, experi\^encias e express\~oes que a pesquisa pretende abarcar;
\item Na constru\c{c}\~ao da amostra, privilegiar sujeitos que detenham os atributos relevantes para o tema de pesquisa;
\item Definir grupos priorit\'arios para o tema de pesquisa, no caso de investiga\c{c}\~oes com amostras heterog\^eneas e abrangentes;
\item Examinar n\~ao apenas as experi\^encias dos indiv\'iduos e grupos, mas como esses grupos interagem entre si;
\item Trabalhar em uma perspectiva de inclus\~ao progressiva das descobertas de pesquisa, confrontando, aos poucos, as evid\^encias coletadas;
\item Ter aten\c{c}\~ao a informa\c{c}\~oes espec\'ificas e que n\~ao se repetem, visto o potencial explicativo e contribui\c{c}\~ao para o tema de pesquisa;
\item Considerar um n\'umero suficiente de pessoas interlocutoras, garantindo complementaridade entre os conte\'udos das entrevistas;
\item Observar se o quadro emp\'irico da pesquisa est\'a plenamente mapeado e bem compreendido;
\item Implementar triangula\c{c}\~ao das t\'ecnicas e m\'etodos de pesquisa.
\end{itemize}

Al\'em disso, para a sele\c{c}\~ao dos participantes, a amostragem n\~ao probabil\'istica pode mobilizar as seguintes estrat\'egias, seguindo a sistematiza\c{c}\~ao realizada por \citet{rocha2020teoria} e \citet{campos2022amostragem}:

\begin{itemize}
\item \textit{Por conveni\^encia}: este tipo de estrat\'egia se baseia na viabilidade da pesquisa, considerando a proximidade ou disponibilidade das pessoas participantes. A sele\c{c}\~ao por conveni\^encia \'e especialmente \'util quando os contatos s\~ao de dif\'icil acesso, ou quando a pesquisa se encontra em est\'agio explorat\'orio.
\item \textit{Intencional}: na sele\c{c}\~ao intencional de pessoas participantes, a equipe pesquisadora define quais pessoas interlocutoras possuem conhecimento significativo sobre o tema pesquisado, bem como suas experi\^encias est\~ao pr\'oximas do objeto de estudo.
\item \textit{Bola de neve}: na sele\c{c}\~ao por bola de neve, as pessoas participantes da pesquisa indicam novas pessoas para participar, que seguem com novas indica\c{c}\~oes. Assim, cria-se uma esp\'ecie de rede de participantes, que compartilham uma experi\^encia ou um entendimento sobre um assunto.
\item \textit{Contatos intersticiais}: nesse \'ultimo caso, as pessoas participantes se encontram em uma mesma situa\c{c}\~ao cotidiana, como em uma fila de espera em um servi\c{c}o de sa\'ude, e assim, podem compartilhar informa\c{c}\~oes a partir do contato da pessoa pesquisadora.
\item \textit{Cotas}: ao possuir conhecimento pr\'evio sobre as caracter\'isticas de uma popula\c{c}\~ao, a pessoa pesquisadora identifica grupos que devem ser contemplados na sele\c{c}\~ao da amostra. Assim, segmentos populacionais espec\'ificos s\~ao garantidos no desenho da pesquisa.
\end{itemize}

\'E pouco usual que estudos qualitativos definam \textit{a priori} o tamanho da amostra. Geralmente, para a determina\c{c}\~ao do n\'umero de pessoas participantes a serem entrevistadas, utiliza-se o chamado \textbf{crit\'erio de satura\c{c}\~ao}, tendo em vista o momento no qual as entrevistas deixam de acrescentar fatores relevantes sobre o objeto de pesquisa. Para o denominado ``fechamento'' da amostra (defini\c{c}\~ao das pessoas participantes a serem inclu\'idas na coleta de dados), \citet{fontanella2011amostragem} prop\~oem a agrega\c{c}\~ao dos trechos de entrevista em enunciados mais simples, e para cada entrevista, s\~ao identificados quantos novos enunciados emergem do conte\'udo. A satura\c{c}\~ao seria alcan\c{c}ada quando novos enunciados deixam de surgir das entrevistas.

O crit\'erio de satura\c{c}\~ao te\'orica tem sido objeto de debate na literatura especializada, em raz\~ao de sua complexa operacionaliza\c{c}\~ao, tanto na defini\c{c}\~ao de um ponto de corte quanto na determina\c{c}\~ao do momento em que, de fato, novas informa\c{c}\~oes deixam de acrescentar elementos relevantes \`a compreens\~ao do fen\^omeno investigado. Uma abordagem alternativa e mais pragm\'atica consiste em reconhecer o ponto em que a continuidade da coleta de dados se torna contraproducente, isto \'e, quando as informa\c{c}\~oes adicionais n\~ao ampliam significativamente a densidade anal\'itica nem contribuem de modo substantivo para o entendimento do objeto de estudo \citep{mason2010sample}.

O n\'umero de pessoas participantes e a diversidade de seus perfis variam amplamente em estudos qualitativos. Dependendo do tema e da estrat\'egia de pesquisa, podem ser delineadas amostras homog\^eneas, compostas por pessoas que compartilham caracter\'isticas e experi\^encias semelhantes, ou amostras heterog\^eneas, marcadas pela diversidade de perfis, contextos e perspectivas. Na sele\c{c}\~ao de pessoas participantes, diferentes t\'ecnicas de amostragem intencional podem ser empregadas para a escolha de casos m\'ultiplos, com o prop\'osito de ampliar a compreens\~ao do objeto de estudo. Essas t\'ecnicas orientam-se, em geral, por dois princ\'ipios fundamentais: o da \textit{diversifica\c{c}\~ao}, que busca captar a heterogeneidade dos casos, e o da \textit{satura\c{c}\~ao}, que visa ao alcance de um n\'ivel adequado de aprofundamento anal\'itico. As principais modalidades de amostragem qualitativa s\~ao apresentadas a seguir, com base nas contribui\c{c}\~oes de \citet{pires2014amostragem}:

\begin{description}
\item[Amostragem por contraste.] A amostragem por contraste tem como objetivo assegurar a inclus\~ao de ao menos uma pessoa participante para cada grupo de interesse, considerado analiticamente relevante para a investiga\c{c}\~ao. Ao selecionar casos que evidenciam diferen\c{c}as significativas entre grupos ou situa\c{c}\~oes, essa abordagem cria uma esp\'ecie de ``mosaico'' de experi\^encias e perspectivas, permitindo ao pesquisador comparar, identificar padr\~oes e compreender varia\c{c}\~oes dentro do fen\^omeno estudado. Trata-se de uma estrat\'egia particularmente \'util em pesquisas qualitativas que buscam explorar contrastes e diversidades contextuais.

\item[Amostragem por homogeneiza\c{c}\~ao.] A amostragem por homogeneiza\c{c}\~ao tem como objetivo concentrar a sele\c{c}\~ao de pessoas participantes em um grupo relativamente homog\^eneo, de modo a explorar em profundidade as semelhan\c{c}as e diferen\c{c}as internas. Assim, s\~ao selecionadas pessoas de um mesmo grupo, analisando-o extensivamente. Ao garantir a diversifica\c{c}\~ao de caracter\'isticas ou experi\^encias dentro desse grupo, a estrat\'egia permite captar a m\'axima variedade de perspectivas intragrupo, oferecendo uma compreens\~ao detalhada e rica do fen\^omeno estudado em contextos semelhantes.

\item[Amostragem por contraste-aprofundamento.] A amostragem por contraste-aprofundamento situa-se entre a abordagem de estudo de caso \'unico e a de m\'ultiplos estudos de caso. Seu objetivo \'e selecionar um pequeno n\'umero de casos espec\'ificos que possam ser analisados em profundidade, permitindo compara\c{c}\~oes e a identifica\c{c}\~ao de paralelos. Dessa forma, a estrat\'egia combina o detalhamento caracter\'istico do estudo de caso \'unico com a possibilidade de identificar padr\~oes e diferen\c{c}as relevantes entre os casos selecionados, enriquecendo a compreens\~ao do fen\^omeno investigado.

\item[Amostragem por contraste-satura\c{c}\~ao.] Esse tipo de amostragem baseia-se na realiza\c{c}\~ao de entrevistas ou na an\'alise de documentos envolvendo m\'ultiplos casos, tratados com menor profundidade que na amostragem por contraste-aprofundamento. O objetivo \'e maximizar a comparabilidade entre os casos, ao mesmo tempo em que se captam diferentes percep\c{c}\~oes e experi\^encias dentro de cada grupo, permitindo identificar padr\~oes, semelhan\c{c}as e varia\c{c}\~oes em uma perspectiva mais ampla.
\end{description}
% ============================================================
% 4. CONSIDERAÇÕES ÉTICAS
% ============================================================
\chapter{Considera\c{c}\~oes \'eticas na coleta de dados}

A coleta de dados prim\'arios envolve pessoas, comunidades e institui\c{c}\~oes --- e isso vincula a equipe de pesquisa a um conjunto de obriga\c{c}\~oes \'eticas que atravessam todo o processo. Como sintetiza o \textit{Guia CLEAR de Monitoramento e Avalia\c{c}\~ao de Pol\'iticas P\'ublicas} \citep{fgvclear2025mea}, a \'etica em avalia\c{c}\~ao n\~ao \'e etapa do trabalho --- \'e dimens\~ao transversal que vai da defini\c{c}\~ao dos objetivos \`a divulga\c{c}\~ao dos resultados.

Este cap\'itulo trata da \'etica em tr\^es frentes. A primeira s\~ao os princ\'ipios gerais e o marco normativo, brasileiro e internacional. A segunda s\~ao as decis\~oes \'eticas que envolvem a sele\c{c}\~ao das participantes --- consentimento informado, equidade e prote\c{c}\~ao de grupos vulner\'aveis. A terceira \'e a \'etica da intera\c{c}\~ao durante a coleta e do tratamento dos dados, que retorna em maior detalhe nas se\c{c}\~oes seguintes do Guia.

\section{Princ\'ipios gerais e marco normativo}

A primeira coisa a reconhecer numa coleta \'etica \'e que toda decis\~ao metodol\'ogica tem implica\c{c}\~oes sociais. Quem participa de uma pesquisa o faz por confian\c{c}a --- de que ser\'a tratado com respeito, de que suas informa\c{c}\~oes ser\~ao protegidas, de que a pesquisa tem prop\'osito leg\'itimo. Romper essa confian\c{c}a compromete n\~ao s\'o o estudo em curso, mas a credibilidade futura de toda a comunidade de avalia\c{c}\~ao.

No Brasil, o tratamento de dados pessoais em pesquisa \'e regulado por dois instrumentos principais. A \textit{Lei Geral de Prote\c{c}\~ao de Dados Pessoais} (LGPD --- Lei n\textsuperscript{o} 13.709/2018) estabelece princ\'ipios para o tratamento de informa\c{c}\~oes pessoais --- finalidade, necessidade, transpar\^encia, seguran\c{c}a e respeito aos direitos da pessoa titular. A \textit{Resolu\c{c}\~ao CNS n\textsuperscript{o} 510/2016}, voltada a pesquisas em Ci\^encias Humanas e Sociais, define os procedimentos de aprecia\c{c}\~ao \'etica de projetos pelo Sistema CEP/CONEP (Comit\^es de \'Etica em Pesquisa e Comiss\~ao Nacional de \'Etica em Pesquisa).

No plano internacional, organismos como o \citet{jpal2020practical} e o \citet{worldbank2021data} t\^em consolidado diretrizes espec\'ificas para pesquisa avaliativa, com \^enfase em transpar\^encia, reprodutibilidade e auditabilidade. N\~ao substituem as exig\^encias brasileiras, mas as complementam com boas pr\'aticas internacionais.

Antes do in\'icio da coleta, projetos que envolvam entrevistas, experimentos ou coleta de informa\c{c}\~oes sens\'iveis devem ser submetidos a um \textbf{Comit\^e de \'Etica em Pesquisa (CEP)}. O parecer atesta que o estudo respeita direitos, privacidade e bem-estar das participantes, com conformidade a padr\~oes nacionais e internacionais.

A \'etica na coleta articula tr\^es frentes: \textbf{(i)} como se selecionam e abordam as participantes; \textbf{(ii)} como se conduz a intera\c{c}\~ao durante a coleta; \textbf{(iii)} como se tratam os dados produzidos. As pr\'oximas subse\c{c}\~oes tratam das duas primeiras e remetem ao Cap\'itulo~6, onde o tratamento dos dados \'e detalhado.

\section{\'Etica na sele\c{c}\~ao de pessoas participantes}

Decidir quem entra na pesquisa \'e, ao mesmo tempo, decis\~ao metodol\'ogica e \'etica. Al\'em de atender a crit\'erios estat\'isticos e de cobertura, o processo de sele\c{c}\~ao precisa assegurar tratamento justo, respeito \`a privacidade e inclus\~ao equitativa. Pesquisas avaliativas envolvem pessoas que podem ser direta ou indiretamente afetadas pelos resultados, o que coloca \`a equipe pesquisadora o dever de proteger seus direitos e minimizar danos \citep{fgvclear2025mea}.

O \textit{FGV CLEAR (2025)} recomenda um conjunto de boas pr\'aticas para a etapa de sele\c{c}\~ao: registrar detalhadamente o processo de sele\c{c}\~ao e/ou sorteio de casos, incluindo listas, c\'odigos, estrat\'egias de acesso e crit\'erios; documentar substitui\c{c}\~oes, perdas e desvios de campo com justificativas verific\'aveis; capacitar as equipes de coleta e supervis\~ao para reconhecer e evitar vieses involunt\'arios; assegurar consentimento informado, explicando \`as participantes os objetivos da pesquisa, o uso dos dados e o direito de recusa; garantir confidencialidade e anonimiza\c{c}\~ao, sobretudo em bases sens\'iveis (educa\c{c}\~ao, sa\'ude, renda, g\^enero); e reconhecer as especificidades de grupos vulner\'aveis, evitando que crit\'erios de amostragem reforcem desigualdades preexistentes ou excluam segmentos marginalizados.

Em avalia\c{c}\~oes experimentais, o \citet{jpal2020practical} acrescenta uma camada espec\'ifica: transpar\^encia e auditabilidade em todas as etapas do processo amostral e de randomiza\c{c}\~ao. Sorteios e designa\c{c}\~oes de grupos de tratamento e controle devem ser documentados de forma reproduz\'ivel --- idealmente com c\'odigos de programa\c{c}\~ao e arquivos p\'ublicos.

\section{\'Etica na intera\c{c}\~ao durante a coleta e no tratamento dos dados}

A \'etica n\~ao termina na sele\c{c}\~ao. Permeia tamb\'em a forma como a equipe se relaciona com as participantes em campo, e como os dados produzidos s\~ao armazenados, analisados e descartados ap\'os o uso autorizado.

\textbf{Na intera\c{c}\~ao durante a coleta}, a postura da equipe entrevistadora \'e determinante. Diferen\c{c}as de idade, g\^enero, posi\c{c}\~ao institucional ou \textit{status} social entre quem entrevista e quem participa podem criar barreiras ou induzir respostas voltadas \`a agradabilidade. Mitigar essas assimetrias exige treinamento \'etico e relacional, escuta ativa, postura neutra e sensibilidade ao verbal e ao n\~ao verbal. Esses cuidados aparecem em detalhe na Se\c{c}\~ao~6.4 (Entrevistas), com as tr\^es etapas do processo: prepara\c{c}\~ao, condu\c{c}\~ao e p\'os-entrevista.

\textbf{No tratamento dos dados}, os princ\'ipios da LGPD orientam coleta, armazenamento, uso e descarte de informa\c{c}\~oes pessoais. Os procedimentos pr\'aticos --- declara\c{c}\~ao de consentimento, mecanismos de anonimiza\c{c}\~ao, controle de acesso, tempo de guarda, elimina\c{c}\~ao segura --- est\~ao na Se\c{c}\~ao~6.3 (Prote\c{c}\~ao dos dados), que articula a fundamenta\c{c}\~ao legal com o lado operacional.

% ============================================================
% 5. ELABORANDO QUESTIONÁRIOS E ROTEIROS
% ============================================================
\chapter{Elaborando question\'arios e roteiros de entrevista}

Question\'arios e roteiros materializam a teoria do estudo nas perguntas que efetivamente chegam a quem responde. \'E a\'i que a coleta ganha ou perde qualidade --- a maioria dos vieses de mensura\c{c}\~ao em uma pesquisa nasce da forma como uma pergunta foi formulada, ordenada ou apresentada. Este cap\'itulo trata dos pontos que mais costumam comprometer essa qualidade.

\section{Formula\c{c}\~ao das quest\~oes em question\'arios}

Esta se\c{c}\~ao percorre os principais aspectos do desenho de question\'arios, da formula\c{c}\~ao das quest\~oes \`a revis\~ao final. Boa parte se aplica tamb\'em a entrevistas, embora as especificidades destas sejam tratadas adiante em t\'opico pr\'oprio.

\subsection{Princ\'ipios b\'asicos de um bom question\'ario}

Segundo \citet{iarossi2006power}, um bom question\'ario deve observar alguns princ\'ipios fundamentais no momento da sua elabora\c{c}\~ao. Abaixo, est\~ao listados os principais conceitos e suas defini\c{c}\~oes:

\begin{itemize}
\item \textbf{Relev\^ancia:} cada pergunta deve ter um prop\'osito claro e estar diretamente ligada aos objetivos da pesquisa. Perguntas ``curiosas'' ou n\~ao essenciais devem ser evitadas, pois aumentam o tempo de resposta sem agregar valor.
\item \textbf{Precis\~ao:} as perguntas devem ser formuladas de maneira a eliciar informa\c{c}\~oes corretas e espec\'ificas, com defini\c{c}\~oes claras de termos, unidades de medida e prazos de refer\^encia.
\item \textbf{Disposi\c{c}\~ao do respondente (\textit{willingness}):} o question\'ario deve considerar a disposi\c{c}\~ao da pessoa participante em colaborar, evitando perguntas excessivamente sens\'iveis, invasivas ou cansativas.
\item \textbf{Clareza:} linguagem simples e acess\'ivel \'e essencial para evitar ambiguidades ou mal-entendidos.
\item \textbf{Consist\^encia:} quest\~oes semelhantes devem seguir o mesmo formato, garantindo comparabilidade.
\item \textbf{Viabilidade:} a extens\~ao do question\'ario deve ser compat\'ivel com o tempo dispon\'ivel para a aplica\c{c}\~ao, respeitando os limites de aten\c{c}\~ao do respondente.
\end{itemize}

\subsection{Escolha do tipo de pergunta}

A escolha do tipo de pergunta \'e uma etapa essencial na elabora\c{c}\~ao de um question\'ario, pois determina tanto a qualidade das informa\c{c}\~oes coletadas quanto a experi\^encia da pessoa respondente. Em linhas gerais, as perguntas podem ser \textbf{abertas ou fechadas}.

\textbf{Perguntas abertas} permitem que a pessoa respondente formule livremente sua resposta, oferecendo maior riqueza de detalhes e nuances. No entanto, elas exigem maior esfor\c{c}o cognitivo e tempo, o que pode gerar desconforto ou mesmo evas\~ao, sobretudo em instrumentos autoadministrados.

J\'a as \textbf{perguntas fechadas} s\~ao mais objetivas e r\'apidas de responder, podendo assumir diferentes formatos: respostas dicot\^omicas (sim/n\~ao, verdadeiro/falso), m\'ultipla escolha, matriz \textit{ranking} ou escalas de avalia\c{c}\~ao (Likert). A escolha do tipo espec\'ifico de pergunta deve ser guiada pelo indicador final que se deseja calcular e pelo m\'etodo de an\'alise. Por exemplo, para calcular uma taxa m\'edia de satisfa\c{c}\~ao, uma escala Likert \'e apropriada; para escolhas \'unicas, \'e mais indicado usar m\'ultipla escolha.

A recomenda\c{c}\~ao mais comum \'e privilegiar perguntas fechadas, pois garantem comparabilidade e reduzem a carga do respondente. Contudo, \'e v\'alido reservar espa\c{c}o para uma ou duas perguntas abertas ao final, ou incluir campos de coment\'ario opcionais. \'E igualmente importante evitar o uso excessivo de diferentes formatos em um mesmo question\'ario, j\'a que isso pode confundir ou desmotivar o participante. Pesquisas sugerem limitar-se a tr\^es tipos de formatos em um mesmo instrumento.

A clareza sobre a finalidade da pergunta ajuda tanto \`a pessoa avaliadora a ser precisa no que deseja medir, quanto \`a pessoa respondente a compreender rapidamente o que se espera dela. Para ajudar nesse processo, \'e interessante agrupar as perguntas em categorias. Segundo \citet{patton2002qualitative}, \'e poss\'ivel definir seis diferentes categorias de quest\~oes:

\begin{itemize}
\item \textbf{Experi\^encia e comportamento:} buscam captar a\c{c}\~oes concretas, experi\^encias ou rotinas --- o que a pessoa faz ou pretende fazer;
\item \textbf{Opini\~oes e valores:} voltadas a compreender processos cognitivos e interpretativos;
\item \textbf{Sentimentos:} buscam rea\c{c}\~oes emocionais a experi\^encias ou situa\c{c}\~oes, normalmente respondidas como ansioso, feliz, com medo, intimidado ou confiante;
\item \textbf{Conhecimento:} investigam informa\c{c}\~oes sobre os conhecimentos da pessoa entrevistada;
\item \textbf{Sensoriais:} relacionadas ao que se v\^e, sente ou ouve;
\item \textbf{Contexto/Demogr\'aficas:} idade, escolaridade, ocupa\c{c}\~ao e outras caracter\'isticas do entrevistado.
\end{itemize}

Al\'em da tipologia, tamb\'em \'e preciso considerar o \textbf{tipo de informa\c{c}\~ao} a ser mensurada \citep{kopper2020survey}. Quando se trata de \textbf{fatos}, recomenda-se utilizar m\'odulos j\'a validados em pesquisas semelhantes, pensar no n\'ivel de precis\~ao necess\'ario e, sempre que poss\'ivel, trabalhar com categorias ou intervalos em vez de estimativas exatas, que podem ser mais dif\'iceis para o respondente. Quest\~oes factuais tamb\'em devem ser adaptadas ao contexto local, garantindo que listas de op\c{c}\~oes reflitam a realidade da popula\c{c}\~ao --- como no caso de categorias de escolaridade ou tipos de ocupa\c{c}\~ao.

As perguntas \textbf{subjetivas}, por sua vez, destinam-se a medir cren\c{c}as, percep\c{c}\~oes e expectativas, como satisfa\c{c}\~ao com servi\c{c}os ou planos para o futuro. Escalas como Likert --- com respostas variando de ``concordo plenamente'' a ``discordo plenamente'' --- ou percentuais s\~ao instrumentos \'uteis, embora seja comum que respostas se concentrem em n\'umeros arredondados. O uso de recursos visuais (como escadas de n\'iveis ou carinhas de express\~ao) pode facilitar a compreens\~ao, especialmente em popula\c{c}\~oes com menor escolaridade. Ainda assim, \'e preciso considerar que perguntas subjetivas podem demandar mais tempo e esfor\c{c}o do respondente, devendo ser usadas com parcim\^onia.

Quando se abordam \textbf{temas sens\'iveis} --- como renda, pr\'aticas ilegais, sa\'ude sexual ou viol\^encia ---, o risco de vi\'es por desejabilidade social aumenta. Para lidar com isso, \'e fundamental assegurar confidencialidade, escolher estrat\'egias apropriadas de coleta e seguir algumas estrat\'egias focadas nesse tipo de quest\~ao. Entre elas, destacam-se:

\begin{itemize}
\item Criar um ambiente confort\'avel e assegurar confidencialidade dos dados para que a pessoa respondente sinta que suas informa\c{c}\~oes est\~ao seguras;
\item Utilizar coleta por question\'arios autoadministrados pode reduzir o medo ou constrangimento do respondente;
\item Evitar que as perguntas sens\'iveis fiquem no in\'icio ou final do question\'ario, pois isso pode evitar que a pessoa respondente se recuse a responder ou se sinta muito desconfort\'avel.
\end{itemize}

Outra categoria a ser considerada s\~ao as \textbf{perguntas que exigem estimativas ou lembran\c{c}a ao longo do tempo}, como gastos, consumo ou renda em per\'iodos extensos. Essas quest\~oes est\~ao mais sujeitas a erros de mem\'oria e baixa precis\~ao. Estrat\'egias \'uteis incluem reduzir o per\'iodo de recorda\c{c}\~ao, utilizar di\'arios ou medi\c{c}\~oes frequentes, pedir respostas em unidades mais familiares ao respondente (como tempo de caminhada em vez de dist\^ancia em quil\^ometros) e inserir verifica\c{c}\~oes de consist\^encia dentro do question\'ario.

Por fim, h\'a tamb\'em as \textbf{perguntas observacionais}, nas quais a pessoa entrevistadora registra diretamente informa\c{c}\~oes observadas no ambiente ou no respondente, como a condi\c{c}\~ao da moradia, a presen\c{c}a de determinados bens ou o estado de conserva\c{c}\~ao de equipamentos. Esse tipo de informa\c{c}\~ao \'e menos sujeito a vieses de resposta, mas pode sofrer com o chamado efeito Hawthorne --- quando as pessoas modificam seu comportamento por estarem sendo observadas. Para reduzir esse risco, \'e importante padronizar protocolos, treinar a equipe entrevistadora e, sempre que poss\'ivel, triangular os dados observacionais com outras fontes.

Assim, a escolha do tipo de pergunta deve ser orientada pelo equil\'ibrio entre a utilidade da informa\c{c}\~ao, a carga cognitiva imposta ao respondente e a qualidade esperada dos dados.

\subsection{\textit{Wording} --- como formular perguntas de forma clara}

A reda\c{c}\~ao das perguntas --- o \textit{wording} --- \'e um dos pontos mais delicados na elabora\c{c}\~ao de um question\'ario. Pequenas mudan\c{c}as na formula\c{c}\~ao podem alterar significativamente as respostas obtidas. Por isso, aten\c{c}\~ao especial deve ser dada \`a escolha das palavras, ao tempo verbal empregado, \`a clareza da linguagem e ao formato das alternativas de resposta.

N\~ao existe uma teoria universalmente aceita sobre formula\c{c}\~ao de quest\~oes, mas h\'a consenso quanto a alguns princ\'ipios b\'asicos que diferenciam boas de m\'as perguntas. De modo geral, uma pergunta deve ser \textbf{Breve, Objetiva, Simples e Espec\'ifica (BOSS)}:

\paragraph{Breve (\textit{Brief}).} Perguntas curtas s\~ao mais f\'aceis de compreender e reduzem a probabilidade de erros de leitura ou interpreta\c{c}\~ao, tanto por parte do entrevistador quanto do respondente. Quest\~oes excessivamente longas costumam incluir mais de uma ideia (perguntas ``escondidas'') ou detalhamento desnecess\'ario, o que aumenta a confus\~ao e dificulta a resposta.

Ser breve, por\'em, n\~ao significa simplesmente encurtar a frase: deseja-se encontrar a forma mais simples de transmitir a ideia, sem sacrificar clareza ou precis\~ao. Em alguns casos, especialmente quando o tema exige que a pessoa entrevistada recupere mem\'orias ou aborde t\'opicos sens\'iveis, perguntas um pouco mais longas podem ser necess\'arias para situar o contexto. No geral, por\'em, a regra \'e dividir assuntos complexos em partes mais simples e evitar incluir perguntas impl\'icitas ou ocultas.

\textit{Orienta\c{c}\~oes pr\'aticas:}
\begin{itemize}
\item Evite perguntas com mais de 20 palavras \citep{payne1951art};
\item Pergunte uma coisa de uma vez e evite perguntas ``escondidas'';
\item Subdivida temas e conceitos complexos em algumas quest\~oes diretas.
\end{itemize}

\paragraph{Objetivo (\textit{Objective}).} Uma boa pergunta deve ser neutra e n\~ao sugerir respostas. Isso significa evitar palavras carregadas de ju\'izo de valor ou que induzam a pessoa entrevistada a concordar ou discordar de forma autom\'atica. Perguntas que come\c{c}am com ``voc\^e n\~ao acha que...'' ou que trazem apenas parte das alternativas poss\'iveis s\~ao exemplos de formula\c{c}\~oes tendenciosas que distorcem os resultados. At\'e mesmo o conjunto de respostas oferecido em uma quest\~ao fechada pode influenciar as escolhas, caso as op\c{c}\~oes n\~ao representem de maneira equilibrada as possibilidades existentes. Por exemplo, se apenas alternativas positivas, como ``bom'', ``muito bom'' e ``excelente'' forem apresentadas, estamos, na pr\'atica, orientando a resposta.

Al\'em disso, termos carregados emocionalmente ou com conota\c{c}\~ao pol\'itica, como ``justo'' ou ``honesto'', devem ser usados com cautela, pois podem ativar associa\c{c}\~oes indesejadas. O ideal \'e manter um tom neutro e informativo, evitando enviesamentos sutis que prejudiquem a qualidade dos dados.

\textit{Orienta\c{c}\~oes pr\'aticas:}
\begin{itemize}
\item Evite perguntas que sugiram respostas (``Voc\^e n\~ao acha que...?'');
\item N\~ao use termos com carga emocional (``justo'', ``honesto'', ``correto'');
\item Ofere\c{c}a alternativas equilibradas e completas nas quest\~oes fechadas.
\end{itemize}

\paragraph{Simples (\textit{Simple}).} Perguntas devem ser formuladas com palavras claras, de uso comum, e sem a presen\c{c}a de jarg\~oes t\'ecnicos, express\~oes regionais ou ambiguidades. O uso de termos t\'ecnicos s\'o deve ser feito quando inevit\'avel e, nesses casos, acompanhado de uma explica\c{c}\~ao breve para garantir que todos os respondentes compartilhem o mesmo entendimento.

Da mesma forma, \'e importante evitar constru\c{c}\~oes gramaticais complexas, como negativas ou duplas negativas, que exigem esfor\c{c}o cognitivo adicional e podem levar a interpreta\c{c}\~oes equivocadas. O objetivo \'e reduzir ao m\'aximo a carga mental necess\'aria para compreender a quest\~ao, sem renunciar \`a precis\~ao. Uma linguagem simples e direta garante que diferentes pessoas, com distintos n\'iveis de escolaridade ou experi\^encia, possam responder de maneira consistente e compar\'avel.

\textit{Orienta\c{c}\~oes pr\'aticas:}
\begin{itemize}
\item Use palavras e express\~oes de conhecimento comum entre todos os participantes;
\item Evite utilizar jarg\~oes t\'ecnicos ou, quando inevit\'avel, explique os conceitos t\'ecnicos;
\item N\~ao use frases com dupla negativa ou que exijam esfor\c{c}o excessivo de interpreta\c{c}\~ao.
\end{itemize}

\paragraph{Espec\'ifico (\textit{Specific}).} Perguntas vagas tendem a gerar respostas igualmente vagas, o que compromete a utilidade das informa\c{c}\~oes coletadas. Termos imprecisos como ``frequentemente'', ``geralmente'' ou ``normalmente'' podem ser interpretados de maneiras muito distintas entre os respondentes, resultando em dados pouco confi\'aveis. O mesmo ocorre com perguntas que abordam mais de um tema ao mesmo tempo --- chamadas de ``duplas'' ou \textit{double-barreled} ---, que confundem e dificultam a resposta.

Al\'em disso, em perguntas fechadas, \'e fundamental garantir que as op\c{c}\~oes de resposta sejam mutuamente exclusivas e abrangentes, de modo que o entrevistado consiga se posicionar de forma adequada. Ser espec\'ifico, portanto, \'e uma quest\~ao de clareza e de respeito \`a capacidade da pessoa respondente de oferecer informa\c{c}\~oes precisas e relevantes, ao mesmo tempo em que aumenta a comparabilidade e a consist\^encia dos dados coletados.

\textit{Orienta\c{c}\~oes pr\'aticas:}
\begin{itemize}
\item Evite termos como ``frequentemente'', ``geralmente'', ``normalmente'', que podem ter significados diferentes para cada pessoa;
\item Evite perguntas duplas (\textit{double-barreled}), que abordam mais de um tema ao mesmo tempo;
\item As alternativas de resposta em quest\~oes fechadas devem ser mutuamente exclusivas e coletivamente exaustivas.
\end{itemize}

\subsection{Sequenciamento}

A ordem das perguntas em um question\'ario \'e t\~ao importante quanto a formula\c{c}\~ao de cada quest\~ao individual. Um bom sequenciamento deve facilitar a administra\c{c}\~ao da entrevista, estimular a colabora\c{c}\~ao da pessoa respondente e reduzir riscos de confus\~ao ou abandono. A seguir, apresentam-se princ\'ipios para organizar a sequ\^encia das perguntas.

\paragraph{Abertura das perguntas.} As perguntas iniciais t\^em papel crucial na cria\c{c}\~ao de confian\c{c}a e engajamento. No in\'icio da entrevista, \'e comum que as pessoas respondentes se sintam inseguras ou desconfiadas quanto ao estudo. Por isso, recomenda-se iniciar com quest\~oes simples, neutras e relacionadas ao tema da pesquisa, permitindo que a pessoa participante se sinta confort\'avel e capaz de responder. Perguntas irrelevantes ou desconectadas do objetivo do levantamento devem ser evitadas, pois podem gerar estranhamento. Por exemplo, em uma pesquisa empresarial, \'e mais apropriado abrir com ``Voc\^e poderia descrever brevemente a sua empresa?'' do que com perguntas sobre h\'abitos pessoais de lazer.

\paragraph{Fluxo das perguntas.} Ap\'os a abertura, as quest\~oes devem acompanhar o racioc\'inio natural da pessoa respondente. Perguntas relacionadas devem ser agrupadas, e mudan\c{c}as de tema devem ocorrer de forma gradual, com transi\c{c}\~oes expl\'icitas como ``Agora, vamos falar sobre...'' ou ``A seguir, algumas perguntas sobre...''. Esse tipo de transi\c{c}\~ao ajuda a manter a organiza\c{c}\~ao e reduz a sensa\c{c}\~ao de descontinuidade. Tamb\'em \'e recomend\'avel que perguntas mais gerais antecedam as espec\'ificas e que o question\'ario seja organizado em blocos tem\'aticos. Al\'em disso, o sequenciamento pode facilitar a recorda\c{c}\~ao quando as perguntas s\~ao organizadas de forma temporal (por exemplo, do per\'iodo mais recente ao mais antigo).

\paragraph{Uso de filtros.} O uso de perguntas filtro \'e fundamental para garantir a relev\^ancia e a efici\^encia do question\'ario. Elas permitem direcionar apenas determinadas se\c{c}\~oes para as pessoas respondentes que possuem experi\^encia ou informa\c{c}\~ao sobre o tema em quest\~ao. Esse recurso evita constrangimentos e respostas for\c{c}adas, como no caso em que se pergunta diretamente a todos sobre ``a taxa de juros do empr\'estimo'' sem antes verificar se a pessoa entrevistada possui um empr\'estimo. Al\'em disso, os filtros ajudam a diferenciar respostas ``n\~ao se aplica'' de ``n\~ao sei'', aumentando a precis\~ao dos dados.

\paragraph{Localiza\c{c}\~ao das perguntas sens\'iveis.} Quest\~oes relacionadas a renda, sa\'ude, sexualidade, comportamentos ilegais ou opini\~oes pol\'iticas tendem a gerar desconforto. A coloca\c{c}\~ao dessas perguntas no in\'icio do question\'ario deve ser evitada, pois pode comprometer a confian\c{c}a da pessoa respondente e levar ao abandono da entrevista. Tamb\'em n\~ao \'e recomend\'avel concentr\'a-las apenas no final, pois h\'a risco de a pessoa interromper sua participa\c{c}\~ao antes de responder a pontos-chave. A pr\'atica mais indicada \'e inseri-las em um momento intermedi\'ario, ap\'os a constru\c{c}\~ao de confian\c{c}a, e dentro de blocos em que fa\c{c}am sentido no fluxo do question\'ario. Al\'em disso, recomenda-se introduzi-las gradualmente, por meio de perguntas de aquecimento que reduzam a sensa\c{c}\~ao de invas\~ao ou julgamento.

\subsection{\textit{Layout} e formata\c{c}\~ao}

O \textit{layout} do question\'ario desempenha um papel fundamental na qualidade dos dados coletados. Um desenho inadequado pode gerar erros de aplica\c{c}\~ao, dificultar a codifica\c{c}\~ao e reduzir a precis\~ao das respostas. Embora muitas vezes negligenciado em favor da economia de espa\c{c}o ou impress\~ao, o \textit{layout} deve priorizar a clareza e a funcionalidade, assegurando que o instrumento seja de f\'acil uso tanto para pessoas entrevistadoras quanto para respondentes.

\paragraph{Identifica\c{c}\~ao e numera\c{c}\~ao.} Cada question\'ario deve conter um identificador \'unico, atribu\'ido previamente, que permita vincular o formul\'ario \`a respectiva pessoa respondente ou unidade amostral. Al\'em disso, as perguntas devem ser numeradas de forma sequencial e cont\'inua em todo o instrumento, mesmo quando o question\'ario estiver dividido em se\c{c}\~oes. Essa pr\'atica evita confus\~oes, facilita o armazenamento e agiliza o processamento dos dados.

\paragraph{Espa\c{c}amento e organiza\c{c}\~ao visual.} O espa\c{c}amento entre perguntas deve ser suficiente para permitir uma aplica\c{c}\~ao clara, al\'em de oferecer espa\c{c}o para anota\c{c}\~oes de pessoas entrevistadoras, editoras ou codificadoras. Tentativas de economizar espa\c{c}o comprimindo perguntas em uma mesma p\'agina podem comprometer a qualidade das informa\c{c}\~oes obtidas. Sempre que poss\'ivel, recomenda-se imprimir em apenas um lado da folha, alinhar caixas de resposta \`a margem e, em casos de perguntas mais longas, recorrer a tabelas que organizem op\c{c}\~oes de forma clara. Quest\~oes abertas tamb\'em devem dispor de espa\c{c}o adequado para o tamanho esperado da resposta.

\paragraph{Instru\c{c}\~oes.} As instru\c{c}\~oes s\~ao parte essencial do question\'ario e precisam estar visivelmente diferenciadas das perguntas. \'E recomend\'avel utilizar caixa alta, negrito, ou coloc\'a-las em caixas de texto, de modo a evitar confus\~oes durante a aplica\c{c}\~ao. Instru\c{c}\~oes devem aparecer imediatamente acima das quest\~oes \`as quais se referem, em vez de serem concentradas no in\'icio do question\'ario ou em manuais separados. \'E fundamental que indiquem claramente como responder (ex.: ``marque apenas uma op\c{c}\~ao'' ou ``assinale todas as op\c{c}\~oes que se aplicam'').

\paragraph{Fonte, formata\c{c}\~ao e \^enfase.} O uso adequado de fontes e estilos contribui para interpreta\c{c}\~oes consistentes. Palavras ou express\~oes-chave devem ser destacadas em negrito ou sublinhado, especialmente quando marcam mudan\c{c}as de tempo ou escopo da pergunta (por exemplo, ``nos \'ultimos 12 meses''). Em question\'arios bil\'ingues, \'e recomend\'avel utilizar fontes distintas para cada idioma, assegurando clareza na leitura.

\paragraph{S\'imbolos e elementos gr\'aficos.} S\'imbolos como c\'irculos, caixas de sele\c{c}\~ao, setas e cores podem ser usados para guiar pessoas entrevistadoras e respondentes, facilitar a navega\c{c}\~ao e indicar fluxos de salto. A padroniza\c{c}\~ao de s\'imbolos aumenta a clareza. Por exemplo, caixas podem indicar locais para n\'umeros, c\'irculos para marca\c{c}\~oes de op\c{c}\~ao e setas para pulos de pergunta. Em pesquisas com mais de uma parte, pode ser \'util usar pap\'eis de cores diferentes para cada se\c{c}\~ao.

\paragraph{Estrutura geral e apresenta\c{c}\~ao profissional.} Um question\'ario bem apresentado deve incluir t\'itulo, data da vers\~ao, informa\c{c}\~oes de contato e, preferencialmente, uma breve introdu\c{c}\~ao que explique o prop\'osito da pesquisa, mesmo quando houver carta de apresenta\c{c}\~ao separada. Essa introdu\c{c}\~ao deve esclarecer a finalidade do estudo, como os dados ser\~ao usados, garantir a confidencialidade das informa\c{c}\~oes e estimar o tempo de preenchimento. Para question\'arios enviados por correio, recomenda-se incluir envelope de devolu\c{c}\~ao pr\'e-pago ou instru\c{c}\~oes claras de retorno.

Em suma, um \textit{layout} bem planejado deve equilibrar clareza, organiza\c{c}\~ao e est\'etica, sem comprometer a precis\~ao dos dados em nome da economia de espa\c{c}o ou papel. O objetivo central \'e oferecer uma ferramenta pr\'atica e profissional, que facilite o trabalho de pessoas entrevistadoras, a compreens\~ao das pessoas respondentes e o processamento subsequente das informa\c{c}\~oes.

\subsection{Tradu\c{c}\~ao}

O processo de tradu\c{c}\~ao \'e fundamental para garantir a equival\^encia conceitual e a clareza do instrumento em diferentes contextos culturais e lingu\'isticos. N\~ao se trata apenas de transpor palavras de um idioma para outro, mas de assegurar que os conceitos centrais mantenham o mesmo significado entre culturas distintas. Um mesmo termo pode ter interpreta\c{c}\~oes divergentes dependendo do pa\'is ou grupo social, e at\'e mesmo quando o conceito \'e equivalente, os indicadores usados para medi-lo podem variar.

A tradu\c{c}\~ao deve buscar a equival\^encia conceitual, e n\~ao apenas literal. Para alcan\c{c}ar esse objetivo, recomenda-se a t\'ecnica de \textit{back translation} \citep{iarossi2006power}. Esse processo envolve quatro passos:

\begin{enumerate}
\item Traduzir o question\'ario do idioma original para o idioma de destino.
\item Traduzir a vers\~ao resultante de volta ao idioma original por um tradutor diferente, ou seja, gerar uma vers\~ao retrotraduzida.
\item Comparar a vers\~ao original com a retrotraduzida, identificando e ajustando discrep\^ancias.
\item Repetir o processo at\'e que n\~ao haja inconsist\^encias relevantes.
\end{enumerate}

\'E essencial que a tradu\c{c}\~ao seja feita por pessoas treinadas, que dominem os dois idiomas e compreendam o conte\'udo tem\'atico do question\'ario. Al\'em disso, recomenda-se utilizar o pr\'e-teste como ferramenta adicional para avaliar a qualidade da tradu\c{c}\~ao, observando se as pessoas respondentes compreendem as quest\~oes conforme o esperado.

Em contextos multil\'ingues, nem sempre \'e poss\'ivel produzir vers\~oes integrais em todas as l\'inguas locais. Nestes casos, parte do instrumento pode ser traduzida em manuais de aplica\c{c}\~ao ou materiais de apoio, devendo-se contar com entrevistadores bil\'ingues ou int\'erpretes capacitados. Em qualquer caso, a linguagem empregada deve ser simples, acess\'ivel e pr\'oxima ao uso cotidiano, evitando registros acad\^emicos ou excessivamente formais.

\section{Boas pr\'aticas na elabora\c{c}\~ao de roteiros de entrevistas}

O objetivo desta se\c{c}\~ao \'e apresentar boas pr\'aticas para a elabora\c{c}\~ao de roteiros de entrevista, considerando tanto o grau de estrutura\c{c}\~ao das perguntas quanto a clareza e adequa\c{c}\~ao da linguagem utilizada. Ser\~ao discutidas, ainda, estrat\'egias para fomentar um ambiente de escuta qualificada, que favore\c{c}a a confian\c{c}a, a abertura e a espontaneidade das pessoas participantes. Embora alguns aspectos relacionados \`a constru\c{c}\~ao de roteiros j\'a tenham sido abordados nos t\'opicos anteriores, especialmente no que se refere \`a elabora\c{c}\~ao de question\'arios, esta se\c{c}\~ao se concentrar\'a em desafios e especificidades pr\'oprias dos estudos qualitativos, oferecendo orienta\c{c}\~oes pr\'aticas para o desenvolvimento de entrevistas.

Os roteiros de entrevista s\~ao instrumentos que orientam a atua\c{c}\~ao das equipes entrevistadoras durante a coleta de dados, contemplando a apresenta\c{c}\~ao da pesquisa e os temas a serem abordados. No processo de elabora\c{c}\~ao desses roteiros, \'e essencial definir de forma clara os objetivos da coleta de dados e explicitar a contribui\c{c}\~ao esperada de cada pergunta para o alcance dos prop\'ositos investigativos. Recomenda-se, ainda, organizar o roteiro em blocos tem\'aticos que re\'unam quest\~oes vinculadas a eixos anal\'iticos comuns do estudo. Essa organiza\c{c}\~ao contribui tanto para sistematizar o trabalho da equipe pesquisadora quanto para favorecer a participa\c{c}\~ao das pessoas entrevistadas, promovendo maior fluidez na transi\c{c}\~ao entre t\'opicos e assegurando a coer\^encia da narrativa produzida.

Na elabora\c{c}\~ao das perguntas, a defini\c{c}\~ao do tipo de entrevista constitui uma etapa fundamental, uma vez que o formato e o grau de detalhamento do roteiro variam de acordo com o n\'ivel de estrutura\c{c}\~ao adotado na coleta de dados \citep{wilson2012research}. Nas entrevistas estruturadas, as perguntas devem ser replicadas a todas as pessoas participantes, a fim de garantir comparabilidade dos resultados. J\'a nas entrevistas semiestruturadas, o roteiro \'e composto por perguntas previamente elaboradas, que servem como guia para assegurar uma base comum de comparabilidade entre as pessoas participantes, preservando, contudo, a flexibilidade necess\'aria para que a pessoa pesquisadora ajuste a ordem, a formula\c{c}\~ao ou o n\'ivel de aprofundamento conforme o curso da intera\c{c}\~ao. Por sua vez, nas entrevistas n\~ao estruturadas, o roteiro tende a restringir-se a t\'opicos amplos ou eixos tem\'aticos gerais, que apenas orientam o di\'alogo, conferindo \`a pessoa entrevistada liberdade para construir e desenvolver suas narrativas, interpreta\c{c}\~oes e significados.

O desenvolvimento de um roteiro de entrevista geralmente inicia-se pela defini\c{c}\~ao do tema de pesquisa, seguida do mapeamento dos t\'opicos e subtemas que dever\~ao ser contemplados. A partir desse levantamento, elaboram-se quest\~oes que explorem aspectos espec\'ificos de cada t\'opico, organizadas de modo progressivo, iniciando por perguntas mais simples e descritivas, que favorecem o engajamento do participante, e avan\c{c}ando gradualmente para quest\~oes mais complexas e anal\'iticas.

Dialogando com as contribui\c{c}\~oes de \citet{eppich2019in}, \'e poss\'ivel listar um conjunto de etapas que podem orientar a constru\c{c}\~ao de roteiros de entrevista:

\begin{itemize}
\item Definir o arcabou\c{c}o te\'orico ou conceitual que orientar\'a o estudo, servindo de base para a identifica\c{c}\~ao dos temas centrais, categorias anal\'iticas e perguntas relevantes \`a investiga\c{c}\~ao.
\item Estabelecer um conjunto de procedimentos operacionais que norteiem a intera\c{c}\~ao entre pessoas pesquisadoras e participantes --- incluindo o n\'ivel de formalidade esperado, orienta\c{c}\~oes sobre o uso da linguagem, postura do entrevistador e etapas da intera\c{c}\~ao.
\item Preparar uma explica\c{c}\~ao breve e acess\'ivel sobre os prop\'ositos da pesquisa, a ser apresentada no in\'icio da entrevista, de modo a garantir transpar\^encia, \'etica e compreens\~ao por parte dos participantes.
\item Identificar as informa\c{c}\~oes sociodemogr\'aficas das pessoas participantes (como sexo, idade, escolaridade, ocupa\c{c}\~ao, entre outras vari\'aveis pertinentes), as quais devem ser registradas durante o processo de entrevista --- preferencialmente no in\'icio da intera\c{c}\~ao, ap\'os o consentimento do participante.
\item Iniciar o roteiro com uma pergunta ampla e introdut\'oria, que funcione como ponto de partida para o di\'alogo, favorecendo o engajamento dos participantes e estimulando o compartilhamento espont\^aneo de suas experi\^encias e percep\c{c}\~oes.
\item Desenvolver perguntas subsequentes alinhadas aos objetivos do estudo e ao perfil das pessoas participantes, organizadas de forma progressiva, de modo a explorar t\'opicos espec\'ificos e incentivar a descri\c{c}\~ao de experi\^encias concretas e exemplos situados, o que contribui para a profundidade e a materialidade da an\'alise.
\item Encerrar a entrevista com perguntas de fechamento, que convidem as pessoas participantes a destacar aspectos que considerem mais relevantes, acrescentar informa\c{c}\~oes n\~ao abordadas anteriormente e refletir sobre os temas discutidos. Esse momento tamb\'em pode ser utilizado para confirmar interpreta\c{c}\~oes preliminares ou impress\~oes da pessoa pesquisadora, fortalecendo a validade dos dados coletados.
\end{itemize}

Antes de iniciar a coleta de dados em escala ampliada, recomenda-se a realiza\c{c}\~ao de um teste piloto, que possibilita avaliar a clareza e a pertin\^encia das perguntas, bem como a fluidez, a sequ\^encia e a dura\c{c}\~ao do roteiro de entrevista. Essa etapa \'e fundamental para identificar poss\'iveis ambiguidades, redund\^ancias ou lacunas, al\'em de verificar o n\'ivel de compreens\~ao, envolvimento e conforto dos participantes diante das quest\~oes propostas. Com base nos resultados do piloto, o roteiro pode ser refinado, contribuindo para o aumento da confiabilidade, consist\^encia e qualidade metodol\'ogica da coleta de dados.

\'E recomendado registrar, no roteiro, instru\c{c}\~oes detalhadas \`a equipe entrevistadora, com orienta\c{c}\~oes sobre como apresentar a pesquisa \`as pessoas participantes, aplicar perguntas condicionais e lidar com situa\c{c}\~oes sens\'iveis durante a intera\c{c}\~ao. Essas instru\c{c}\~oes contribuem para assegurar maior consist\^encia e rigor nos dados coletados. Al\'em disso, vimos que, no in\'icio das entrevistas, cabe solicitar que as pessoas participantes falem de forma ampla sobre seu perfil e trajet\'oria, utilizando perguntas simples e acess\'iveis que favore\c{c}am a cria\c{c}\~ao de um clima de confian\c{c}a e proximidade. Esse momento inicial funciona como uma etapa de aquecimento, permitindo que o entrevistado se sinta mais \`a vontade para compartilhar suas percep\c{c}\~oes e experi\^encias. Ap\'os essa introdu\c{c}\~ao, podem ser apresentadas perguntas mais espec\'ificas e complexas, diretamente relacionadas aos objetivos da pesquisa, de modo a aprofundar gradualmente a explora\c{c}\~ao do tema.

Ao longo da entrevista, \'e fundamental empregar uma linguagem clara e adequada ao perfil do p\'ublico de interesse, evitando jarg\~oes t\'ecnicos ou termos amb\'iguos que possam comprometer a compreens\~ao. Recomenda-se, ainda, observar e, quando pertinente, incorporar express\~oes ou formas de comunica\c{c}\~ao pr\'oprias do grupo entrevistado, pois isso contribui para fortalecer a empatia, a proximidade e a fluidez do di\'alogo. Essa aten\c{c}\~ao \`a linguagem amplia a qualidade da intera\c{c}\~ao e favorece que as pessoas participantes compreendam plenamente as perguntas, expressando-se com autenticidade e seguran\c{c}a. Tamb\'em \'e importante incluir no roteiro quest\~oes de fechamento, que ofere\c{c}am \`as pessoas entrevistadas a oportunidade de recapitular os principais pontos abordados e destacar os aspectos que considerem mais relevantes. Esse momento final pode abrir espa\c{c}o para que acrescentem informa\c{c}\~oes que n\~ao tenham sido contempladas anteriormente.

Nos roteiros de entrevista, as perguntas devem privilegiar o ``como'' dos processos, uma vez que esse tipo de formula\c{c}\~ao tende a ser mais produtivo para captar significados, pr\'aticas e racionalidades subjacentes \`as experi\^encias das pessoas participantes, apresentando, assim, maior potencial explicativo e interpretativo para os prop\'ositos da pesquisa qualitativa \citep{albaret2023semistructured}. Al\'em disso, \'e fundamental que as perguntas sejam claras, concisas e focalizadas, evitando a combina\c{c}\~ao de m\'ultiplos temas em uma mesma formula\c{c}\~ao, o que pode gerar ambiguidades ou dispersar o foco da conversa \citep{driscoll2011introduction}.

Por fim, recomenda-se que a equipe pesquisadora formule perguntas que estimulem a reflex\~ao, a descri\c{c}\~ao de experi\^encias concretas e a explicita\c{c}\~ao de significados, refor\c{c}ando a coer\^encia entre o roteiro e os objetivos anal\'iticos da investiga\c{c}\~ao.

Outro conjunto de recomenda\c{c}\~oes se dirige \`as caracter\'isticas das perguntas em um roteiro de entrevista:

\begin{itemize}
\item \'E fundamental evitar formula\c{c}\~oes que possam ser respondidas apenas com ``sim'' ou ``n\~ao'', pois esse tipo de pergunta tende a restringir o n\'ivel de detalhamento e a riqueza do discurso dos participantes.
\item De maneira geral, cabe \`a equipe entrevistadora estimular as pessoas entrevistadas a descrever processos, experi\^encias e percep\c{c}\~oes com suas pr\'oprias palavras \citep{roulston2013asking}.
\item O roteiro deve evitar perguntas indutivas, que sugiram respostas desejadas ou direcionem a narrativa em determinada dire\c{c}\~ao, comprometendo a espontaneidade e a validade das informa\c{c}\~oes coletadas. \'E prefer\'ivel priorizar perguntas abertas, que promovam a reflex\~ao, a descri\c{c}\~ao de experi\^encias e a express\~ao de significados.
\item O uso de sondagens e quest\~oes de aprofundamento pode auxiliar o entrevistador a incentivar a amplia\c{c}\~ao das respostas, mantendo a neutralidade e evitando interfer\^encias indevidas no conte\'udo.
\end{itemize}

Em s\'intese, a elabora\c{c}\~ao de roteiros de entrevista constitui uma etapa central no desenho metodol\'ogico da pesquisa qualitativa, pois traduz, em termos operacionais, as escolhas te\'oricas e anal\'iticas do estudo. Al\'em disso um roteiro bem estruturado favorece a constru\c{c}\~ao de um ambiente de escuta qualificada, no qual as pessoas participantes possam expressar livremente suas experi\^encias e significados.

% ============================================================
% 6. APLICAÇÃO DOS QUESTIONÁRIOS
% ============================================================
\chapter{Aplica\c{c}\~ao dos question\'arios}

A forma como os question\'arios s\~ao aplicados tem impacto direto sobre a qualidade e a confiabilidade dos dados. Cada m\'etodo de coleta --- seja presencial, telef\^onico, online ou h\'ibrido --- exige cuidados espec\'ificos com a abordagem dos participantes, a log\'istica de campo, os aspectos \'eticos e o treinamento das equipes envolvidas.

Como destaca o \textit{Guia CLEAR de Monitoramento e Avalia\c{c}\~ao de Pol\'iticas P\'ublicas} \citep{fgvclear2025mea}, o momento da aplica\c{c}\~ao deve ser entendido como uma etapa estrat\'egica do processo avaliativo, que transforma o planejamento metodol\'ogico em pr\'atica concreta.

Neste cap\'itulo, s\~ao apresentados os principais elementos a considerar durante a aplica\c{c}\~ao dos question\'arios, incluindo: \textbf{(i)} a realiza\c{c}\~ao de testes-piloto e o aprimoramento dos instrumentos; \textbf{(ii)} estrat\'egias para aumentar a taxa de resposta; \textbf{(iii)} procedimentos de prote\c{c}\~ao e confidencialidade dos dados; \textbf{(iv)} boas pr\'aticas para a condu\c{c}\~ao de entrevistas; e \textbf{(v)} mecanismos de controle de qualidade durante e ap\'os a coleta.

\section{Piloto}

Os \textbf{testes-piloto} s\~ao uma etapa fundamental entre o desenho do question\'ario e sua aplica\c{c}\~ao definitiva. Segundo o \textit{Guia CLEAR de Avalia\c{c}\~ao de Impacto} \citep{fgvclear2025ai}, essa fase tem como principal objetivo antecipar e corrigir eventuais problemas de compreens\~ao das perguntas, estrutura, fluxo ou dura\c{c}\~ao do instrumento.

O piloto deve reproduzir, o m\'aximo poss\'ivel, as condi\c{c}\~oes reais de campo, mesmo que em escala reduzida. Isso permite avaliar se as perguntas s\~ao claras, se sua sequ\^encia faz sentido e se o tempo total da entrevista \'e adequado. Tamb\'em \'e o momento de testar ferramentas e dispositivos eletr\^onicos (como \textit{tablets}, formul\'arios digitais e \textit{softwares} de coleta) e de ajustar o treinamento da equipe, garantindo que todas as pessoas integrantes compreendam as instru\c{c}\~oes e sigam os mesmos procedimentos.

Ao final, recomenda-se elaborar um relat\'orio de retroalimenta\c{c}\~ao, registrando as dificuldades encontradas, as sugest\~oes da equipe entrevistadora e as percep\c{c}\~oes das pessoas respondentes. Essas informa\c{c}\~oes devem orientar ajustes de linguagem, ordem, formato de resposta e instru\c{c}\~oes de campo \citep{iarossi2006power, oecd2022data, fgvclear2025mea}.

\section{Aumentando a taxa de resposta e evitando n\~ao respostas}

A \textbf{taxa de resposta} \'e um dos principais indicadores da qualidade de uma pesquisa. Quando o n\'umero de pessoas participantes efetivas \'e muito baixo, cresce o risco de \textbf{vi\'es de n\~ao resposta}, o que pode comprometer a representatividade da amostra e distorcer as conclus\~oes \citep{duflo2007using, jpal2020practical}. Esse tipo de vi\'es \'e mais comum em contextos de vulnerabilidade social, em coletas remotas com baixa conectividade ou em pesquisas sobre temas sens\'iveis --- por exemplo, sa\'ude mental, renda ou desempenho escolar.

Altas taxas de n\~ao resposta podem ocorrer por tr\^es caminhos principais: a recusa em participar, o abandono no meio da aplica\c{c}\~ao ou as falhas de registro. Cada um desses caminhos exige estrat\'egias espec\'ificas de preven\c{c}\~ao e mitiga\c{c}\~ao. A boa pr\'atica come\c{c}a por reconhecer que a n\~ao resposta n\~ao \'e apenas um problema operacional, mas um indicador de qualidade do trabalho de campo que deve ser monitorado sistematicamente.

Para minimizar o problema, \'e fundamental adotar estrat\'egias de engajamento adequadas ao perfil do p\'ublico e ao tipo de aplica\c{c}\~ao. Em levantamentos presenciais, o foco deve estar na postura e na comunica\c{c}\~ao da equipe de campo; em coletas online, na clareza da mensagem e na experi\^encia da pessoa usu\'aria. Entre as pr\'aticas mais eficazes, destacam-se:

\begin{itemize}
\item \textbf{Garantir anonimato e confidencialidade}, explicando claramente como as informa\c{c}\~oes ser\~ao utilizadas e protegidas;
\item \textbf{Apresentar o prop\'osito da pesquisa de forma transparente}, mostrando que a participa\c{c}\~ao contribui para o aprimoramento de pol\'iticas e programas p\'ublicos;
\item \textbf{Realizar contato pr\'evio e comunica\c{c}\~ao clara} sobre os objetivos e benef\'icios da pesquisa, preparando as pessoas participantes para a coleta;
\item \textbf{Enviar lembretes e convites personalizados}, especialmente em coletas eletr\^onicas --- estudos apontam que mensagens curtas e diretas podem elevar a taxa de resposta em at\'e 20\% \citep{oecd2022data};
\item \textbf{Manter o question\'ario conciso e relevante}, evitando repeti\c{c}\~oes e perguntas irrelevantes que cansam o respondente;
\item \textbf{Oferecer incentivos simb\'olicos ou financeiros}, quando adequado ao contexto;
\item \textbf{Treinar a equipe entrevistadora} para conduzir a aplica\c{c}\~ao com empatia, clareza e respeito, criando um ambiente de confian\c{c}a e colabora\c{c}\~ao.
\end{itemize}

Quando as perdas amostrais ultrapassam limites aceit\'aveis --- a literatura geralmente indica como refer\^encia inicial perdas acima de 10\% --- recomenda-se considerar \textbf{reposi\c{c}\~ao de casos e amostragem complementar}, sempre com \textbf{documenta\c{c}\~ao transparente} dos motivos de evas\~ao e dos crit\'erios usados para substituir observa\c{c}\~oes perdidas.

Segundo a \citet{oecd2022data}, o registro sistem\'atico das taxas de n\~ao resposta constitui um indicador-chave de qualidade de campo em pesquisas, devendo ser incorporado aos relat\'orios t\'ecnicos. Esse registro permite identificar padr\~oes de aus\^encia ou recusa, avaliar potenciais vieses de sele\c{c}\~ao e ajustar estrat\'egias de coleta de dados --- contribuindo para a robustez, confiabilidade e transpar\^encia dos resultados. Essas medidas, combinadas de forma coerente com o desenho metodol\'ogico, ajudam a melhorar a qualidade dos dados e a reduzir custos com recontatos e substitui\c{c}\~oes durante o campo.

O quadro-resumo abaixo apresenta as principais estrat\'egias para elevar a taxa de resposta.

\begin{table}[H]
\centering
\caption{Estrat\'egias pr\'aticas para aumentar a taxa de resposta}
\label{tab:taxa_resposta}
\begin{tabular}{p{0.30\linewidth} p{0.32\linewidth} p{0.30\linewidth}}
\toprule
\textbf{Estrat\'egia} & \textbf{Aplicabilidade} & \textbf{Benef\'icio principal} \\
\midrule
Garantia de anonimato & Pesquisas sens\'iveis (ex.: satisfa\c{c}\~ao, sa\'ude) & Reduz receio de exposi\c{c}\~ao \\
Comunica\c{c}\~ao clara & Coletas com p\'ublico amplo & Aumenta engajamento e legitimidade \\
Lembretes e \textit{follow-ups} & Coletas online e pain\'eis & Recupera pessoas participantes inativas \\
Incentivos & Amostras volunt\'arias & Melhora motiva\c{c}\~ao e participa\c{c}\~ao \\
Treinamento da equipe & Entrevistas presenciais & Melhora qualidade da abordagem \\
\bottomrule
\end{tabular}
\figmeta{Elabora\c{c}\~ao pr\'opria.}{}{}
\end{table}

\section{Prote\c{c}\~ao dos dados}

A prote\c{c}\~ao e o uso respons\'avel dos dados constituem um pilar \'etico essencial em qualquer processo de coleta para al\'em das exig\^encias legais vigentes. No Brasil, a \textit{Lei Geral de Prote\c{c}\~ao de Dados Pessoais} (LGPD --- Lei n\textsuperscript{o} 13.709/2018) estabelece princ\'ipios fundamentais para o tratamento de informa\c{c}\~oes pessoais, como \textbf{finalidade, necessidade, consentimento e seguran\c{c}a}.

Conforme orientam o \textit{Guia CLEAR de Monitoramento e Avalia\c{c}\~ao de Pol\'iticas P\'ublicas} \citep{fgvclear2025mea} e as diretrizes da \citet{oecd2022data}, toda pesquisa deve garantir que os dados sejam coletados, armazenados e utilizados de forma \textbf{transparente, proporcional e segura}. Isso inclui:

\begin{itemize}
\item obter consentimento livre e esclarecido de todos os participantes;
\item limitar a coleta \`as informa\c{c}\~oes estritamente necess\'arias para os objetivos da pesquisa;
\item aplicar mecanismos de anonimiza\c{c}\~ao e controle de acesso durante o armazenamento;
\item assegurar a elimina\c{c}\~ao segura dos dados ap\'os o t\'ermino do uso autorizado.
\end{itemize}

Antes do in\'icio da coleta, \'e recomend\'avel que o projeto seja submetido a um Comit\^e de \'Etica em Pesquisa (CEP) --- especialmente quando envolve entrevistas, experimentos ou coleta de informa\c{c}\~oes sens\'iveis. O parecer de aprova\c{c}\~ao \'etica atesta que o estudo respeita os direitos, a privacidade e o bem-estar das pessoas participantes, al\'em de garantir que os procedimentos de consentimento, armazenamento e uso dos dados estejam em conformidade com padr\~oes nacionais e internacionais. Essa exig\^encia \'e alinhada \`as pr\'aticas da \textit{Resolu\c{c}\~ao CNS n\textsuperscript{o} 510/2016} (voltada a pesquisas em Ci\^encias Humanas e Sociais) e \`as orienta\c{c}\~oes internacionais do \citet{worldbank2021data} e do \citet{jpal2020practical}.

Os question\'arios e formul\'arios devem conter uma \textbf{declara\c{c}\~ao clara de consentimento}, informando:

\begin{itemize}
\item o objetivo da pesquisa e o uso previsto dos dados;
\item o per\'iodo de guarda e as medidas de seguran\c{c}a adotadas;
\item o direito da pessoa participante de recusar ou retirar o consentimento a qualquer momento, sem qualquer preju\'izo.
\end{itemize}

Esses cuidados --- aliados \`a revis\~ao \'etica e ao cumprimento da LGPD --- refor\c{c}am a credibilidade institucional e a integridade cient\'ifica da pesquisa.

\section{Entrevistas}

As \textbf{entrevistas} exigem cuidados espec\'ificos para garantir a qualidade das informa\c{c}\~oes e a \'etica da intera\c{c}\~ao. De acordo com o \textit{Guia CLEAR de Avalia\c{c}\~ao de Impacto} \citep{fgvclear2025ai} e \citet{creswell2014research}, o papel da equipe entrevistadora vai al\'em da aplica\c{c}\~ao mec\^anica do roteiro: ela atua como mediadora entre o protocolo da pesquisa e a experi\^encia vivida pela pessoa respondente. Sua postura, preparo e sensibilidade influenciam diretamente a profundidade e a fidedignidade das respostas.

Um aspecto central a ser considerado \'e a \textbf{din\^amica de poder na entrevista}. Diferen\c{c}as de idade, g\^enero, posi\c{c}\~ao institucional ou \textit{status} social entre quem entrevista e quem participa podem criar barreiras de comunica\c{c}\~ao ou induzir respostas voltadas \`a agradabilidade --- fen\^omeno relacionado ao vi\'es de desejabilidade social. Para mitigar essas assimetrias, o \textit{Guia CLEAR de Avalia\c{c}\~ao de Impacto} \citep{fgvclear2025ai} e \citet{creswell2014research} recomendam um \textbf{treinamento \'etico e relacional robusto}, com \^enfase em tr\^es dimens\~oes: a import\^ancia da escuta ativa e da postura neutra; o uso de roteiros de apoio e simula\c{c}\~oes de entrevistas durante o treinamento, capacitando a equipe entrevistadora a conduzir as entrevistas de forma adequada ao contexto; e, sempre que poss\'ivel, a designa\c{c}\~ao de entrevistadores que compartilhem caracter\'isticas culturais ou lingu\'isticas com o p\'ublico-alvo.

Como destacam \citet{eppich2019in}, \`a medida que as pessoas participantes descrevem suas experi\^encias, a equipe entrevistadora atenta identifica quest\~oes relevantes que emergem da conversa e as explora em maior detalhe, relacionando-as com as perguntas de pesquisa do estudo. Para assegurar a qualidade desse processo, \'e recomend\'avel estruturar a entrevista em tr\^es etapas principais: a prepara\c{c}\~ao, a condu\c{c}\~ao e o p\'os-entrevista.

\paragraph{Prepara\c{c}\~ao e pr\'e-entrevista.} Antes da sess\~ao, \'e fundamental:

\begin{itemize}
\item \textbf{Agendar e informar:} combinar data, hor\'ario e local adequados com a pessoa participante, compartilhando antecipadamente informa\c{c}\~oes gerais sobre a pesquisa para garantir um consentimento verdadeiramente informado \citep{guazi2021diretrizes}.
\item \textbf{Preparar o ambiente:} garantir um local confort\'avel, privativo e livre de interrup\c{c}\~oes, seja presencial ou virtual.
\item \textbf{Para entrevistas remotas}, testar previamente a conex\~ao e o funcionamento dos equipamentos.
\end{itemize}

\paragraph{Condu\c{c}\~ao da entrevista.} Durante a intera\c{c}\~ao, a pessoa entrevistadora deve adotar, quando poss\'ivel, as seguintes pr\'aticas, baseadas na sistematiza\c{c}\~ao proposta em \citet{rocha2020teoria}:

\begin{itemize}
\item \textbf{Iniciar com agradecimento e contextualiza\c{c}\~ao:} agradecer a participa\c{c}\~ao, apresentar-se e relembrar o objetivo da pesquisa de forma generalista para evitar vieses.
\item \textbf{Confirmar o consentimento:} no caso de grava\c{c}\~ao, informar e solicitar a autoriza\c{c}\~ao expl\'icita do participante.
\item \textbf{Estabelecer \textit{rapport}:} criar um clima de confian\c{c}a e di\'alogo aberto desde o in\'icio.
\item \textbf{Guiar a conversa estrategicamente:} seguir uma sequ\^encia que parta de perguntas gerais e simples para as mais espec\'ificas e complexas, utilizando linguagem clara e acess\'ivel.
\item \textbf{Explorar aprofundamentos:} realizar perguntas de \textit{follow-up} em momentos-chave para elucidar aspectos relevantes que emergem na conversa.
\item \textbf{Manter uma postura neutra e respeitosa:} ouvir ativamente, evitar julgamentos e observar atentamente as express\~oes n\~ao verbais, que tamb\'em carregam significados importantes.
\item \textbf{Encerrar de forma aberta:} concluir agradecendo novamente a contribui\c{c}\~ao, perguntar se o participante gostaria de acrescentar algo e explicar brevemente como as informa\c{c}\~oes ser\~ao utilizadas.
\end{itemize}

\paragraph{P\'os-entrevista.} Ap\'os a coleta, \'e crucial:

\begin{itemize}
\item \textbf{Transcrever os dados:} realizar a transcri\c{c}\~ao, equilibrando fidedignidade ao discurso original e legibilidade, de acordo com os objetivos do estudo.
\item \textbf{Armazenar com seguran\c{c}a:} garantir o armazenamento \'etico e seguro dos dados (\'audio, v\'ideo, transcri\c{c}\~oes), em estrita conformidade com os protocolos de confidencialidade e prote\c{c}\~ao de dados estabelecidos.
\end{itemize}

Esses cuidados constroem uma rela\c{c}\~ao \'etica e respeitosa entre pesquisadora e participante, e da\'i v\^em dados mais confi\'aveis e interpreta\c{c}\~oes mais ricas.

\section{Controle de qualidade dos dados}

O controle de qualidade precisa ser planejado desde o desenho do instrumento, com mecanismos de verifica\c{c}\~ao cont\'inua que peguem erros de coleta, inconsist\^encias e desvios de protocolo enquanto ainda d\'a tempo de corrigir \citep{fgvclear2025mea,gertler2018impact}. Quatro procedimentos complementares fazem a maior parte do trabalho:

\paragraph{1. Reentrevistas (\textit{back-checks}).} Consistem na replica\c{c}\~ao de parte do question\'ario a uma amostra reduzida dos entrevistados originais (geralmente 5\% a 10\% dos casos). Seu objetivo \'e verificar a consist\^encia das respostas e a atua\c{c}\~ao da equipe de campo, detectando poss\'iveis falhas de registro, erros de digita\c{c}\~ao ou desvios de conduta. Esse procedimento tamb\'em serve como mecanismo de supervis\~ao e de retroalimenta\c{c}\~ao imediata, permitindo corrigir problemas ainda durante o trabalho de campo.

\textit{Exemplo:} reaplicar 10\% dos question\'arios ap\'os tr\^es dias, comparando as respostas originais e as reentrevistas para medir a taxa de consist\^encia.

\paragraph{2. Perguntas de consist\^encia interna.} S\~ao quest\~oes redundantes ou cruzadas, inclu\'idas propositalmente no question\'ario para avaliar a coer\^encia entre respostas relacionadas. Elas ajudam a identificar padr\~oes de resposta autom\'atica, lapsos de aten\c{c}\~ao e inconsist\^encias l\'ogicas. Quando diverg\^encias significativas s\~ao detectadas, o instrumento pode ser ajustado para vers\~oes futuras, fortalecendo a confiabilidade geral da base de dados.

\textit{Exemplo:} comparar a renda declarada com a ocupa\c{c}\~ao ou o n\'umero de filhos com a idade relatada.

\paragraph{3. Supervis\~ao de campo.} A supervis\~ao direta \'e um componente essencial do controle de qualidade operacional. Supervisores devem acompanhar periodicamente as entrevistas, observar o comportamento da equipe entrevistadora e garantir o cumprimento dos protocolos \'eticos e t\'ecnicos.

\textit{Exemplo:} \textit{checklists} di\'arios, relat\'orios de visitas e reuni\~oes de \textit{feedback} ajudam a manter a padroniza\c{c}\~ao e o desempenho das equipes. Essa pr\'atica tamb\'em permite identificar gargalos log\'isticos e refor\c{c}ar o treinamento quando necess\'ario.

\paragraph{4. Auditoria eletr\^onica.} Em coletas realizadas com \textit{tablets} ou formul\'arios digitais, \'e poss\'ivel realizar auditorias eletr\^onicas para detectar padr\~oes an\^omalos nas bases. O uso de tecnologias de \textit{Computer Assisted Personal Interviewing} (CAPI) permite monitorar tempos de resposta, coordenadas GPS, datas e hor\'arios de entrevistas, al\'em de respostas duplicadas ou incompletas. Esses indicadores ajudam a identificar entrevistas fict\'icias (\textit{curbstoning}), erros de entrada ou comportamento inadequado da equipe de campo.

O monitoramento cont\'inuo \'e recomendado por institui\c{c}\~oes como o \citet{jpal2020practical} e o \citet{worldbank2021data}, e por uma raz\~ao pr\'atica: pegar falhas durante o campo \'e barato; pegar depois, na limpeza e valida\c{c}\~ao dos dados, sai muito caro.

% ============================================================
% 7. CONSIDERAÇÕES FINAIS
% ============================================================
\chapter{Considera\c{c}\~oes finais}

\section{Limites da generaliza\c{c}\~ao e custos de opera\c{c}\~ao}

Duas dimens\~oes atravessam qualquer coleta prim\'aria: at\'e onde os achados generalizam e quanto a opera\c{c}\~ao custa. Frequentemente s\~ao tratadas como obst\'aculos, mas s\~ao caracter\'isticas estruturais --- e reconhec\^e-las explicitamente fortalece a credibilidade da pesquisa.

\paragraph{Limita\c{c}\~oes \`a generaliza\c{c}\~ao e validade externa.} Nem toda coleta visa a generalizar estatisticamente. Pesquisas qualitativas, estudos de caso e experimentos locais entregam \textit{insights} profundos, mas \textbf{contextualmente delimitados}. A obriga\c{c}\~ao da equipe \'e explicitar essas limita\c{c}\~oes de validade externa --- relatar contexto, amostra, crit\'erios de sele\c{c}\~ao e poss\'iveis vieses de composi\c{c}\~ao ---, n\~ao tentar minimiz\'a-las. Como observam \citet{gertler2018impact} e o \textit{Guia CLEAR de Avalia\c{c}\~ao de Impacto} \citep{fgvclear2025ai}, a \textbf{transpar\^encia metodol\'ogica} \'e a melhor resposta a limites de generaliza\c{c}\~ao: aumenta tanto a confian\c{c}a da pessoa leitora quanto a utilidade pr\'atica dos resultados.

\paragraph{Custos, log\'istica e restri\c{c}\~oes operacionais.} A log\'istica de coletas de grande escala gera atrasos, sobrecustos e lacunas amostrais. A coleta presencial \'e a mais precisa, mas demanda tempo, deslocamentos, equipe treinada e recursos financeiros expressivos. Os m\'etodos remotos --- telefone, formul\'arios online --- esbarram em barreiras tecnol\'ogicas, baixa ades\~ao e maior variabilidade na qualidade das respostas.

Mitiga-se esses riscos com cronogramas realistas (com margem para imprevistos), com tecnologias de coleta assistida (CAPI/CATI e plataformas digitais para entrevistas remotas) quando o or\c{c}amento aperta, com monitoramento sistem\'atico de indicadores de campo (dura\c{c}\~ao das entrevistas, taxas de resposta, tempos de transi\c{c}\~ao) que permitem corre\c{c}\~oes de rota, e com diversifica\c{c}\~ao dos modos de aplica\c{c}\~ao --- combinando canais para n\~ao deixar segmentos relevantes da popula\c{c}\~ao de fora.

O \textit{Guia CLEAR de Monitoramento e Avalia\c{c}\~ao} \citep{fgvclear2025mea} trata o planejamento log\'istico como componente da estrat\'egia metodol\'ogica, n\~ao como detalhe operacional. Cronograma, aloca\c{c}\~ao de recursos, sele\c{c}\~ao de entrevistadores e defini\c{c}\~ao de locais de coleta determinam tanto a efici\^encia quanto os vieses de campo.

A Tabela~\ref{tab:desafios_mitigacao} consolida os principais desafios discutidos ao longo deste Guia e suas estrat\'egias de mitiga\c{c}\~ao.

\begin{table}[H]
\centering
\caption{Principais desafios na coleta de dados e estrat\'egias de mitiga\c{c}\~ao}
\label{tab:desafios_mitigacao}
\small
\begin{tabular}{p{0.16\linewidth} p{0.20\linewidth} p{0.20\linewidth} p{0.25\linewidth}}
\toprule
\textbf{Desafio} & \textbf{Descri\c{c}\~ao} & \textbf{Efeitos potenciais} & \textbf{Estrat\'egias de mitiga\c{c}\~ao} \\
\midrule
Vi\'es de desejabilidade social & Respondentes ajustam respostas para parecer mais ``adequados'' & Compromete validade interna e distorce padr\~oes & Anonimato, perguntas neutras, treinamento \'etico \\
Formula\c{c}\~ao inadequada & Linguagem t\'ecnica ou amb\'igua dificulta entendimento & Erros de mensura\c{c}\~ao e perda de comparabilidade & Piloto cognitivo, simplifica\c{c}\~ao, consist\^encia de escalas \\
N\~ao resposta e evas\~ao & Recusa, formul\'arios incompletos, perdas & Vi\'es de sele\c{c}\~ao e redu\c{c}\~ao de representatividade & Comunica\c{c}\~ao pr\'evia, lembretes, reposi\c{c}\~ao, documenta\c{c}\~ao \\
Assimetrias de poder & Din\^amica interpessoal afeta autenticidade das respostas & Vi\'es de intera\c{c}\~ao e distor\c{c}\~ao qualitativa & Treinamento \'etico e comportamental, escuta ativa, neutralidade \\
Limita\c{c}\~oes \`a generaliza\c{c}\~ao & Resultados restritos a contextos espec\'ificos & Reduz validade externa e aplicabilidade & Transpar\^encia, triangula\c{c}\~ao, contextualiza\c{c}\~ao \\
Custos e log\'istica & Restri\c{c}\~oes or\c{c}ament\'arias e operacionais & Atrasos, lacunas amostrais, perda de qualidade & Planejamento antecipado, CAPI/CATI, monitoramento de campo \\
\bottomrule
\end{tabular}
\figmeta{Elabora\c{c}\~ao pr\'opria.}{}{}
\end{table}

\section{S\'intese}

A coleta prim\'aria acessa o que registros administrativos e bases secund\'arias n\~ao acessam: mecanismos de implementa\c{c}\~ao, desafios enfrentados pelos atores em campo, modos pelos quais o contexto modula os resultados de uma pol\'itica. Esse \'e o seu valor para a gest\~ao p\'ublica baseada em evid\^encias --- e tamb\'em a raz\~ao pela qual exige rigor: dados prim\'arios mal coletados produzem diagn\'osticos errados, e diagn\'osticos errados produzem recomenda\c{c}\~oes piores do que n\~ao recomendar nada.

Este Guia percorreu as decis\~oes que sustentam uma coleta de qualidade --- m\'etodos quantitativos e qualitativos, crit\'erios de amostragem, considera\c{c}\~oes \'eticas, elabora\c{c}\~ao de instrumentos, cuidados operacionais de campo. Cada uma dessas decis\~oes afeta validade, confiabilidade e utilidade do que se produz. E quase nenhuma sai perfeita de primeira: o que distingue uma coleta robusta de uma coleta apenas formalmente correta n\~ao \'e a aus\^encia de problemas, mas o que se faz com eles --- o reconhecimento pr\'evio dos limites, a documenta\c{c}\~ao das decis\~oes, a flexibilidade de ajustar procedimentos durante o campo.

\'E o que se espera ter compartilhado aqui: uma refer\^encia pr\'atica para quem desenha e conduz pesquisas e avalia\c{c}\~oes no Brasil.

% ============================================================
% REFERÊNCIAS
% ============================================================
\begin{thebibliography}{99}

\bibitem[Albaret e Deas(2023)]{albaret2023semistructured}
ALBARET, M.; DEAS, J. \textbf{Semistructured interviews}. In: BADACHE, F.; KIMBER, L.; MERTENS, L. (orgs.). \textit{International organizations and research methods: an introduction}. Ann Arbor: University of Michigan Press, 2023. p. 82--89.

\bibitem[Almeida(2016)]{almeida2016grupos}
ALMEIDA, R. \textbf{Roteiro para o emprego de grupos focais}. In: ABDAL, A. et al. (orgs.). \textit{M\'etodos de pesquisa em Ci\^encias Sociais: Bloco Qualitativo}. S\~ao Paulo: CEBRAP/SESC, 2016. p. 42--59.

\bibitem[Alonso(2016)]{alonso2016metodos}
ALONSO, A. \textbf{M\'etodos qualitativos de pesquisa: uma introdu\c{c}\~ao}. In: ABDAL, A. et al. (orgs.). \textit{M\'etodos de pesquisa em Ci\^encias Sociais: Bloco Qualitativo}. S\~ao Paulo: CEBRAP/SESC, 2016. p. 8--23.

\bibitem[Bauer e Aarts(2013)]{bauer2013construcao}
BAUER, M.; AARTS, B. \textbf{A constru\c{c}\~ao do corpus: um princ\'ipio para a coleta de dados qualitativos}. In: BAUER, M.; GASKELL, G. (orgs.). \textit{Pesquisa qualitativa com texto, imagem e som}. Petr\'opolis: Vozes, 2013. p. 39--63.

\bibitem[Brasil(2018)]{brasil2018lgpd}
BRASIL. \textbf{Lei n\textsuperscript{o} 13.709/2018 --- Lei Geral de Prote\c{c}\~ao de Dados Pessoais (LGPD)}. Bras\'ilia: Presid\^encia da Rep\'ublica, 2018.

\bibitem[Bryman(2012)]{bryman2012social}
BRYMAN, A. \textbf{Social Research Methods}. 4. ed. Oxford: Oxford University Press, 2012.

\bibitem[Campos e Saidel(2022)]{campos2022amostragem}
CAMPOS, C.; SAIDEL, M. \textbf{Amostragem em investiga\c{c}\~oes qualitativas: conceitos e aplica\c{c}\~oes ao campo da sa\'ude}. \textit{Revista Pesquisa Qualitativa}, v. 10, n. 25, p. 404--424, 2022.

\bibitem[Creswell(2014)]{creswell2014research}
CRESWELL, J. W. \textbf{Research Design: Qualitative, Quantitative, and Mixed Methods Approaches}. 4. ed. Thousand Oaks: SAGE, 2014.

\bibitem[De Paula et al.(2014)]{depaula2014modos}
DE PAULA, C. et al. \textbf{Modos de condu\c{c}\~ao da entrevista em pesquisa fenomenol\'ogica: relato de experi\^encia}. \textit{Revista Brasileira de Enfermagem}, v. 67, n. 3, p. 468--472, 2014.

\bibitem[Driscoll(2011)]{driscoll2011introduction}
DRISCOLL, D. \textbf{Introduction to primary research: observations, surveys, and interviews}. In: LOWE, C.; ZEMLIANSKY, P. (orgs.). \textit{Writing spaces: readings on writing}, v. 2. South Carolina: Parlor Press, 2011. p. 153--174.

\bibitem[Duflo et al.(2007)]{duflo2007using}
DUFLO, E.; GLENNERSTER, R.; KREMER, M. \textbf{Using Randomization in Development Economics Research: A Toolkit}. In: \textit{Handbook of Development Economics}, v. 4. Amsterdam: Elsevier, 2007. p. 3895--3962.

\bibitem[Duru et al.(2021)]{duru2021grant}
DURU, M.; KOPPER, S.; TURITTO, J. \textbf{Grant and budget management}. J-PAL, 2021. Dispon\'ivel em: \url{https://www.povertyactionlab.org/resource/grant-and-budget-management}.

\bibitem[Eppich et al.(2019)]{eppich2019in}
EPPICH, W.; GORMLEY, G.; TEUNISSEN, P. \textbf{In-depth interviews}. \textit{Healthcare Simulation Research}, 2019. DOI: 10.1007/978-3-030-26837-4\_12.

\bibitem[FGV CLEAR(2025a)]{fgvclear2025mea}
FGV CLEAR. \textbf{Guia CLEAR de Monitoramento e Avalia\c{c}\~ao de Pol\'iticas P\'ublicas: do diagn\'ostico \`a decis\~ao}. Rio de Janeiro: FGV CLEAR, 2025.

\bibitem[FGV CLEAR(2025b)]{fgvclear2025ai}
FGV CLEAR. \textbf{Guia CLEAR de Avalia\c{c}\~ao de Impacto}. Rio de Janeiro: FGV CLEAR, 2025.

\bibitem[FGV CLEAR(2025c)]{fgvclear2025poder}
FGV CLEAR. \textbf{An\'alise de Poder Estat\'istico}. Rio de Janeiro: FGV CLEAR, 2025.

\bibitem[Fontanella et al.(2011)]{fontanella2011amostragem}
FONTANELLA, B. et al. \textbf{Amostragem em pesquisas qualitativas: proposta de procedimentos para constatar a satura\c{c}\~ao te\'orica}. \textit{Caderno de Sa\'ude P\'ublica}, v. 27, n. 2, p. 389--394, 2011.

\bibitem[Fraser e Gondim(2004)]{fraser2004fala}
FRASER, M.; GONDIM, S. \textbf{Da fala do outro ao texto negociado: discuss\~oes sobre a entrevista na pesquisa qualitativa}. \textit{Paid\'eia}, v. 14, n. 28, p. 139--152, 2004.

\bibitem[Gertler et al.(2018)]{gertler2018impact}
GERTLER, P. J. et al. \textbf{Impact Evaluation in Practice}. 2. ed. Washington, DC: World Bank, 2018.

\bibitem[Guazi(2021)]{guazi2021diretrizes}
GUAZI, T. \textbf{Diretrizes para o uso de entrevistas semiestruturadas em investiga\c{c}\~oes cient\'ificas}. \textit{Revista Educa\c{c}\~ao, Pesquisa e Inclus\~ao}, v. 2, p. 1--20, 2021.

\bibitem[Iarossi(2006)]{iarossi2006power}
IAROSSI, G. \textbf{The Power of Survey Design}. Washington, DC: World Bank, 2006.

\bibitem[Jim\'enez(2012)]{jimenez2012entrevista}
JIM\'ENEZ, I. \textbf{La entrevista en la investigaci\'on cualitativa: nuevas tendencias y retos}. \textit{Revista Calidad en la Educaci\'on Superior}, v. 3, n. 1, p. 119--139, 2012.

\bibitem[J-PAL(2020)]{jpal2020practical}
J-PAL. \textbf{Practical Guide to Field Research and Randomized Evaluations}. Cambridge, MA: MIT, 2020.

\bibitem[Kabir(2016)]{kabir2016methods}
KABIR, S. M. S. \textbf{Methods of Data Collection}. In: \textit{Basic Guidelines for Research: An Introductory Approach for All Disciplines}. 1. ed. 2016. p. 201--275.

\bibitem[Kopper e Kannan(2020)]{kopper2020survey}
KOPPER, S.; KANNAN, H. \textbf{Survey Designs}. J-PAL, 2020. Dispon\'ivel em: \url{https://www.povertyactionlab.org/blog/3-20-20/best-practices-conducting-phone-surveys}.

\bibitem[Kopper e Parry(2021)]{kopper2021survey}
KOPPER, S.; PARRY, K. \textbf{Survey Design}. J-PAL, 2021. Dispon\'ivel em: \url{https://www.povertyactionlab.org/resource/survey-design}.

\bibitem[Leit\~ao(2024)]{leitao2024entrevista}
LEIT\~AO, C. \textbf{A entrevista como instrumento de pesquisa cient\'ifica: planejamento, execu\c{c}\~ao e an\'alise}. In: PIMENTEL, M.; SANTOS, E. (orgs.). \textit{Metodologia de pesquisa cient\'ifica em inform\'atica na Educa\c{c}\~ao: abordagem qualitativa}. Porto Alegre: SBC, 2024.

\bibitem[Lima(2016)]{lima2016entrevistas}
LIMA, M. \textbf{O uso de entrevistas na pesquisa emp\'irica}. In: ABDAL, A. et al. (orgs.). \textit{M\'etodos de pesquisa em Ci\^encias Sociais: Bloco Qualitativo}. S\~ao Paulo: CEBRAP/SESC, 2016. p. 24--41.

\bibitem[Lotta e Pires(2020)]{lotta2020categorizando}
LOTTA, G.; PIRES, R. \textbf{Categorizando usu\'arios ``f\'aceis'' e ``dif\'iceis'': pr\'aticas cotidianas de implementa\c{c}\~ao de pol\'iticas p\'ublicas e a produ\c{c}\~ao de diferen\c{c}as sociais}. \textit{Dados}, v. 63, n. 4, p. 1--40, 2020.

\bibitem[Mason(2010)]{mason2010sample}
MASON, M. \textbf{Sample size and saturation in PhD studies using qualitative interviews}. \textit{FQS}, v. 11, n. 3, 2010.

\bibitem[Mattos(2011)]{mattos2011resultados}
MATTOS, P. \textbf{``Os resultados desta pesquisa (qualitativa) n\~ao podem ser generalizados'': pondo os pingos nos is de tal ressalva}. \textit{Cadernos EBAPE.BR}, v. 9, p. 450--468, 2011.

\bibitem[Minayo(2017)]{minayo2017amostragem}
MINAYO, M. \textbf{Amostragem e satura\c{c}\~ao em pesquisa qualitativa: consensos e controv\'ersias}. \textit{Revista Pesquisa Qualitativa}, v. 5, n. 7, p. 1--12, 2017.

\bibitem[Morra-Imas e Rist(2009)]{morra2009road}
MORRA-IMAS, L. G.; RIST, R. C. \textbf{The road to results: Designing and conducting effective development evaluations}. Washington, DC: World Bank Publications, 2009.

\bibitem[Neuman(2012)]{neuman2012designing}
NEUMAN, W. L. \textbf{Designing the face-to-face survey}. In: \textit{Handbook of survey methodology for the social sciences}. New York: Springer, 2012. p. 227--248.

\bibitem[Neuman(2014)]{neuman2014social}
NEUMAN, W. L. \textbf{Social Research Methods: Qualitative and Quantitative Approaches}. 7. ed. Pearson Education Limited, 2014.

\bibitem[OECD(2022)]{oecd2022data}
OECD. \textbf{Data Collection and Management for Evaluation and Policy Analysis}. Paris: OECD, 2022.

\bibitem[Patton(1990)]{patton1990qualitative}
PATTON, M. \textbf{Qualitative evaluation and research methods}. California: Sage, 1990.

\bibitem[Patton(2002)]{patton2002qualitative}
PATTON, M. \textbf{Qualitative research \& evaluation methods}. 3. ed. Thousand Oaks: Sage, 2002.

\bibitem[Payne(1951)]{payne1951art}
PAYNE, S. L. \textbf{The Art of Asking Questions}. Princeton: Princeton University Press, 1951.

\bibitem[Pires(2014)]{pires2014amostragem}
PIRES, \'A. P. \textbf{Amostragem e pesquisa qualitativa: ensaio te\'orico e metodol\'ogico}. In: POUPART, J. et al. \textit{A pesquisa qualitativa: enfoques epistemol\'ogicos e metodol\'ogicos}. Petr\'opolis: Vozes, 2014.

\bibitem[Rocha(2020)]{rocha2020teoria}
ROCHA, V. \textbf{Da teoria \`a an\'alise: uma introdu\c{c}\~ao ao uso de entrevistas individuais semiestruturadas na ci\^encia pol\'itica}. \textit{Revista Pol\'itica Hoje}, 2020. DOI: 10.51359/1808-8708.2021.247229.

\bibitem[Roulston(2013)]{roulston2013asking}
ROULSTON, K. \textbf{Asking questions and individual interviews}. Sage Research Methods, 2013.

\bibitem[Silva et al.(2006)]{silva2006entrevista}
SILVA, G. et al. \textbf{Entrevista como t\'ecnica de pesquisa qualitativa}. \textit{Online Brazilian Journal of Nursing}, v. 5, n. 2, p. 246--257, 2006.

\bibitem[Silva e Fortunato(2021)]{silva2021modera}
SILVA, P.; FORTUNATO, M. \textbf{Modera, observa, escuta e foca-te na conversa\c{c}\~ao de grupo: uma reflex\~ao cr\'itica}. In: S\'A, P.; COSTA, A.; MOREIRA, A. (orgs.). \textit{Reflex\~oes em torno de metodologias de investiga\c{c}\~ao: recolha de dados}. Aveiro: UA Editora, 2021. p. 37--51.

\bibitem[Silva e Russo(2019)]{silva2019aplicacao}
SILVA, L.; RUSSO, R. \textbf{Aplica\c{c}\~ao de entrevistas em pesquisas qualitativas}. \textit{Revista de Gest\~ao e Projetos}, v. 10, n. 1, p. 1--6, 2019.

\bibitem[Taherdoost(2021)]{taherdoost2021data}
TAHERDOOST, H. \textbf{Data collection methods and tools for research: a step-by-step guide to choose data collection technique for academic and business research projects}. \textit{International Journal of Academic Research in Management}, v. 10, n. 1, p. 10--38, 2021.

\bibitem[Vaessen et al.(2020)]{vaessen2020evaluation}
VAESSEN, J.; LEMIRE, S.; BEFANI, B. \textbf{Evaluation of International Development Interventions: An Overview of Approaches and Methods}. Washington, DC: World Bank, 2020.

\bibitem[Wilson(2012)]{wilson2012research}
WILSON, V. \textbf{Research methods: interviews}. \textit{Evidence Based Library and Information Practice}, v. 7, p. 96--98, 2012.

\bibitem[World Bank(2021)]{worldbank2021data}
WORLD BANK. \textbf{Data Collection in Fragile States: Innovations from Africa and Beyond}. Washington, DC: World Bank, 2021.

\bibitem[Lima et~al.(2025)]{Lima2025}
Lima, L. e colaboradores (2025). \textit{Guia CLEAR de Monitoramento e Avaliação de Políticas Públicas: do diagnóstico à decisão}. São Paulo: FGV CLEAR.

\end{thebibliography}

\end{document}
